\chapter*{About this blueprint}

This blueprint reproduces M.~Stoll, \emph{Galois groups over $\QQ$ of some iterated
polynomials}, Arch.\ Math.\ \textbf{59} (1992), 239--244, and links each of its definitions and
statements to the Lean~4 formalization in this repository. Chapters~1--3 correspond to the three
sections of the paper; their numbering has been chosen so that the numbered items carry the
\emph{same} numbers as in the paper (Facts~1.0, Lemma~1.1, \dots, Lemma~2.2). Chapter~0 collects
the definitions from the introduction, and Chapter~4 the general-purpose theory that the
formalization develops on the side, in Mathlib generality, under
\texttt{QuadraticIterates/Mathlib/}.

All statements of Chapters~0--4 are formalized and proved: there is no \texttt{sorry} anywhere,
so every node of those chapters is green. Chapter~5 is different in kind: it is the \emph{plan}
for extending the formalization to the results of Li \cite{li2020, li2021}, and none of its
nodes is formalized yet. Where the formal statement differs from the printed one --- a weaker hypothesis,
a total function with a junk value, a more general coefficient ring --- this is recorded in a
\emph{Formalization} paragraph. The Lean names attached to each item are the declarations that
carry its content; clicking one leads to the generated documentation.

\emph{Standing hypotheses.} Throughout, $a$ is an integer such that $-a$ is not a square in
$\QQ$, $f := X^2 + a \in \ZZ[X]$, and $f_0 := X$, $f_{n+1} := f(f_n) = f_n^2 + a$ are its
iterates. $K_n$ is the splitting field of $f_n$ over $\QQ$ and $\Om_n := \Gal(K_n/\QQ) =
\Gal(f_n/\QQ)$. $\Ct{n}$ denotes the $n$-fold iterated wreath product of the cyclic group $C_2$
with itself. In Chapter~2, $g \in \ZZ[X^2]$ is an even polynomial and $\varepsilon = \pm 1$.

% Introduction: the objects of the paper. This chapter is unnumbered (\chapter*), so its items
% are numbered 0.1, 0.2, ... and the paper's own numbering starts in the next chapter.
\chapter*{Introduction}
\addcontentsline{toc}{chapter}{Introduction}

Throughout this paper, $a$ denotes an integer such that $-a$ is not a square in $\QQ$,
$f := X^2 + a \in \ZZ[X]$, and $f_0 := X$, $f_{n+1} := f(f_n) = f_n^2 + a$ for all $n \ge 0$, are
the iterates of $f$. Let $c_1 := -a$ and $c_{n+1} := f(c_n) = c_n^2 + a$
$(= f_n(-a) = f_n(a) = f_{n+1}(0))$ for $n \ge 1$. There exists an integer sequence $(b_n)_{n
\ge 1}$ such that for all $n \ge 1$ we have $c_n = \prod_{d \dvd n} b_d$
(Lemma~\ref{lem:1.1}). Let $K_n$ be the splitting field of $f_n$ over $\QQ$ and denote by
$\Om_n := \Gal(K_n/\QQ) = \Gal(f_n/\QQ)$ its Galois group over the rational numbers.

This paper deals with the problem of determining $\Om_n$. Following R.~W.~K.~Odoni's paper
\cite{odoni}, we generalize his result concerning the case $a = 1$ to get for all $n \ge 1$:
\[ \Om_n \cong \Ct{n}, \text{ if and only if } b_1, b_2, \dots, b_n
   \text{ are $2$-independent in } \QQ \tag{$*$} \]
(Theorem~\ref{thm:section1}). $\Om_n$ always injects into $\Ct{n}$, the $n$-fold wreath product
of $C_2$ with itself. In order to prove the condition on the right hand side of $(*)$ (at least
for some $a \in \ZZ$), we consider more generally iteration sequences associated to certain even
polynomials $g \in \ZZ[X^2]$ (Chapter~2). Specializing these results, we get with $(*)$: if
($a > 0$ and $a \equiv 1$ or $2 \bmod 4$) or ($a < 0$ and $a \equiv 0 \bmod 4$), then
$\Om_n \cong \Ct{n}$ for all $n \ge 1$ (Theorem~\ref{thm:section3}). Moreover, this gives a proof
of the conjecture of John~E.~Cremona \cite{cremona}, who verified numerically that, for $a = 1$,
$|b_n|$ is not a square when $2 \le n \le 5 \cdot 10^7$.

\begin{definition}[The iterates $f_n$]
  \label{def:iterates}
  \lean{QuadraticIterates.iteratedPoly, QuadraticIterates.iteratedPoly_zero,
    QuadraticIterates.iteratedPoly_succ, QuadraticIterates.iteratedPoly_succ_comp,
    QuadraticIterates.iteratedPoly_add, QuadraticIterates.map_iteratedPoly,
    QuadraticIterates.monic_iteratedPoly, QuadraticIterates.natDegree_iteratedPoly,
    QuadraticIterates.ncard_rootSet_iteratedPoly_le}
  \leanok
  $f_0 := X$ and $f_{n+1} := f_n^2 + a$. These are monic of degree $2^n$ --- so $f_n$ has at
  most $2^n$ roots --- and they satisfy
  $f_{m+n} = f_m \circ f_n$; in particular $f_{n+1} = f_n \circ (X^2 + a)$, which is the form in
  which the iterate is taken apart in Chapter~1.
  \formalnote{\texttt{iteratedPoly} is defined over an arbitrary commutative semiring; the
  sequence of the paper is the case $R = \ZZ$. Since forming iterates commutes with ring
  homomorphisms (\lname{map\_iteratedPoly}), the image of $f_n$ in $\QQ[X]$ --- written
  \texttt{f}$\QQ$\texttt{[a, n]} in Lean --- is again an iterate, and all statements over $\QQ$
  are about that polynomial.}
\end{definition}

\begin{definition}[The fields $K_n$ and the groups $\Om_n$]
  \label{def:galois}
  \uses{def:iterates}
  \lean{QuadraticIterates.splittingField, QuadraticIterates.GaloisGroup,
    QuadraticIterates.splittingField.instIsSplittingField,
    QuadraticIterates.splittingField.instFiniteDimensional,
    QuadraticIterates.splittingField.instIsGalois,
    QuadraticIterates.splittingField_zero_eq_bot,
    QuadraticIterates.card_galoisGroup_eq_finrank}
  \leanok
  $K_n \subseteq \Qbar$ is the subfield generated by the roots of $f_n$, and
  $\Om_n = \Gal(f_n/\QQ)$. One has $K_0 = \QQ$ and $\#\Om_n = [K_n : \QQ]$.
  \formalnote{$K_n$ is realized concretely as an \texttt{IntermediateField} of
  \texttt{AlgebraicClosure} $\QQ$ --- namely \texttt{IntermediateField.adjoin} of the root set ---
  rather than as an abstract splitting field, so that the tower $K_n \subseteq K_{n+1}$ is an
  inclusion of subfields of one fixed field and relative degrees are
  \texttt{IntermediateField.relfinrank}. That this field \emph{is} a splitting field, is finite
  over $\QQ$ and is Galois are instances. $\Om_n$ is Mathlib's \texttt{Polynomial.Gal}.}
\end{definition}

\begin{definition}[The iterated wreath power $\Ct{n}$]
  \label{def:wreath}
  \lean{QuadraticIterates.WreathPower, QuadraticIterates.card_wreathPower}
  \leanok
  $\Ct{n}$ is the $n$-fold iterated regular wreath product of $C_2$ with itself; it has order
  $2^{2^n - 1}$.
  \formalnote{This is Mathlib's \texttt{IteratedWreathProduct}, applied to the group
  \texttt{Multiplicative (ZMod 2)}. Statements of the form ``$\Om_n \cong \Ct{n}$'' are
  formalized as \texttt{Nonempty} of a \texttt{MulEquiv} between \texttt{GaloisGroup a n} and
  \texttt{WreathPower n}: the \emph{existence} of an isomorphism, which is what the paper
  asserts, rather than a chosen one.}
\end{definition}

\begin{definition}[The sequence $c_n$]
  \label{def:cseq}
  \uses{def:iterates}
  \lean{QuadraticIterates.cSeq, QuadraticIterates.cSeq_zero, QuadraticIterates.cSeq_one,
    QuadraticIterates.cSeq_succ, QuadraticIterates.cSeq_two,
    QuadraticIterates.cSeq_succ_eq_neg_one_pow_mul_eval,
    QuadraticIterates.intCast_cSeq_eq_ite_mul_eval_zero}
  \leanok
  $c_1 = -a$ and $c_{n+1} = c_n^2 + a$ for $n \ge 1$. One has $c_{n+1} = (-1)^{2^n} f_n(a)$, so
  $c_{n+1} = f_n(a) = f_{n+1}(0)$ for $n \ge 1$, while $c_1 = -f_1(0)$.
  \formalnote{\texttt{cSeq} is total, with the junk value $c_0 = 0$; the recursion
  \lname{cSeq\_succ} therefore carries the hypothesis $n \ge 1$. The junk value is not
  arbitrary: it is what makes $c$ a \emph{strong divisibility sequence}
  ($\gcd(c_m, c_n) = |c_{\gcd(m,n)}|$, see Lemma~\ref{lem:1.1}), which is how integrality of the
  $b_n$ is obtained. \texttt{cSeq} is the specialization to $g = X^2 + a$, $\varepsilon = -1$ of
  the general sequence of Definition~\ref{def:gammaseq}.}
\end{definition}

\begin{definition}[The Möbius factors $b_n$]
  \label{def:bseq}
  \uses{def:cseq}
  \lean{QuadraticIterates.bSeq, QuadraticIterates.bSeq_one,
    QuadraticIterates.bSeq_eq_moebiusFactorR}
  \leanok
  $b_n := \prod_{d \dvd n} c_d^{\mu(n/d)}$, where $\mu$ is the Möbius function. In particular
  $b_1 = c_1 = -a$.
  \formalnote{The paper \emph{constructs} the $b_n$ inside Lemma~1.1~b) and proves integrality
  afterwards. The formalization instead \emph{defines} $b_n$ at once, as an element of $\ZZ$: it
  is \texttt{moebiusFactorR} of the sequence $c$, the preimage in $\ZZ$ of the rational Möbius
  product (junk when that product is not integral). Lemma~\ref{lem:1.1}~b) is then the statement
  that this preimage really maps to the rational product $\prod_{ed = n} c_d^{\mu(e)}$
  (\lname{intCast\_bSeq}), i.e. that the product is an integer. The rational product itself is
  \texttt{moebiusFactorK} with $K = \QQ$, written over the divisor antidiagonal of $n$
  (Lemma~\ref{lem:moebiusfactor}).}
\end{definition}

\begin{definition}[$2$-independence]
  \label{def:twoindep}
  \lean{QuadraticIterates.TwoIndependent, QuadraticIterates.TwoIndependent.ne_zero,
    QuadraticIterates.twoIndependent_iff_linearIndependent,
    QuadraticIterates.TwoIndependent.of_snoc, sqClass}
  \leanok
  Nonzero rational numbers $a_1, \dots, a_n$ are called \emph{$2$-independent} if their residue
  classes in the $\FF_2$-vector space $\QQ^*/(\QQ^*)^2$ are linearly independent.
  \formalnote{\texttt{TwoIndependent} is stated in the equivalent elementary form ``no nonempty
  subfamily has a square product'', which is what the proofs use, for a family of elements of any
  commutative monoid; that all $a_i$ are nonzero is then a consequence, $0$ being a square
  (\lname{TwoIndependent.ne\_zero}). The translation into linear independence of the classes
  $\texttt{sqClass}(a_i)$ in $L^*/(L^*)^2$, for any field $L$, is
  \lname{twoIndependent\_iff\_linearIndependent}. The families of the paper are indexed by
  \texttt{Fin n}, and $c_1, \dots, c_n$ appears as the family $i \mapsto c_{i+1}$.}
\end{definition}

\chapter{\texorpdfstring{A criterion for $\Om_n \cong \Ct{n}$}{A criterion for the full wreath
  power}}

% The counter is set so that the items below carry the numbers 1.0--1.6 of the paper.
\setcounter{theorem}{-1}

In this section, we derive criterion $(*)$ from the introduction. To begin with, we list a few
simple

\begin{facts}[Kummer tower]
  \label{facts:1.0}
  \uses{def:iterates, def:galois}
  \lean{QuadraticIterates.natDegree_iteratedPoly, QuadraticIterates.kummer_degree_bound,
    QuadraticIterates.splittingField_le_succ, QuadraticIterates.splittingField_mono,
    QuadraticIterates.relfinrank_succ_le,
    QuadraticIterates.relfinrank_succ_eq_finrank_adjoin,
    QuadraticIterates.mem_rootSet_iteratedPoly_succ, QuadraticIterates.exists_sq_eq_sub}
  \leanok
  $f_n$ has degree $2^n$. $K_{n+1} = K_n\bigl(\sqrt{\alpha - a} \bigm| \alpha \text{ root of }
  f_n\bigr)$, so $K_n \subset K_{n+1}$ is a $2$-Kummer extension, and
  $[K_{n+1} : K_n] \le 2^{2^n}$.
\end{facts}
\begin{proof}
  \leanok
  \uses{lem:multiquadratic}
  The degree statement is immediate from $f_{n+1} = f_n^2 + a$. A root of $f_{n+1}$ is precisely
  an element $\beta$ with $\beta^2 + a$ a root of $f_n$; as $\Qbar$ is algebraically closed,
  every root $\alpha$ of $f_n$ has such a $\beta$, namely a square root of $\alpha - a$. Hence
  $K_{n+1}$ is generated over $K_n$ by square roots of the $2^n$ elements $\alpha - a \in K_n$,
  and a multiquadratic extension generated by $r$ square roots has degree at most $2^r$
  (Lemma~\ref{lem:multiquadratic}).
\end{proof}

\begin{lemma}[cf.\ \cite{odoni}, Lemma 4.4]
  \label{lem:1.1}
  \uses{def:cseq, def:bseq}
  \lean{QuadraticIterates.cSeq_pos, QuadraticIterates.abs_le_cSeq,
    QuadraticIterates.cSeq_gcd, QuadraticIterates.cSeq_associated_gcd,
    QuadraticIterates.intCast_bSeq, QuadraticIterates.cSeq_eq_prod_bSeq,
    QuadraticIterates.isCoprime_bSeq, QuadraticIterates.bSeq_ne_zero,
    QuadraticIterates.cSeq_ne_zero, QuadraticIterates.cSeq_factorization_shape}
  \leanok
  \begin{enumerate}
    \item[a)] $\forall n \ge 2 : c_n > 0$.
    \item[b)] The $b_n$ $(n \ge 1)$ are integers such that
      $\forall n \ge 1 : c_n = \prod_{d \dvd n} b_d$, and the $b_n$ are pairwise coprime.
  \end{enumerate}
\end{lemma}
\begin{proof}
  \leanok
  \uses{lem:strongdvd, lem:moebiusfactor}
  a) $a \ne 0, -1$ since $-a$ is not a square. Therefore $c_2 = a^2 + a \ge |a| > 0$, and
  induction shows that $c_{n+1} = c_n^2 + a \ge a^2 + a \ge |a| > 0$.

  b) Since all $c_n \ne 0$, the rational number $\prod_{d \dvd n} c_d^{\mu(n/d)}$ is defined, and
  we have to show that it is already an integer. Let $p$ be a prime dividing at least one of the
  $c_n$. Let $m := \min\{n \ge 1 \mid p \dvd c_n\}$ and $e := v_p(c_m)$. Then we have
  $p \dvd c_n$ iff $p^e \dvd c_n$ iff $m \dvd n$. Indeed, with $x_n := f_n(0)$ $(= \pm c_n$ for
  $n \ge 1)$ we have $x_m \equiv 0 \bmod p$ and $\bmod\ p^e$, but $\not\equiv 0 \bmod p^{e+1}$;
  therefore $(x_n)_{n \ge 0}$ is purely periodic $\bmod\ p$ and $\bmod\ p^{e+1}$ with period $m$,
  and mod $p^{e+1}$ we have $x_{m+1} = x_m^2 + a \equiv a \equiv x_1$, so $(x_n)$ is periodic
  from $n = 1$ on with period $m$, hence $x_n$ is never zero mod $p^{e+1}$ if $m \nmid n$. That
  is, $v_p(c_n) = e$ if $m \dvd n$, and $0$ otherwise. Now
  $v_p(b_n) = \sum_{d \dvd n} v_p(c_d)\,\mu(n/d) = e \sum_{d \dvd n/m} \mu(n/(dm)) = e$ if
  $n = m$, and $0$ otherwise. Therefore $b_n$ is integral at $p$, and moreover $p$ divides only
  one of the $b_n$. Hence all the $b_n$ are integers and pairwise coprime; Möbius inversion gives
  $c_n = \prod_{d \dvd n} b_d$.
  \formalnote{The formalization splits this argument into its two halves and proves each in the
  generality it deserves. Integrality (and Möbius inversion) come from \emph{strong
  divisibility}: $\gcd(c_m, c_n) = |c_{\gcd(m,n)}|$ (Lemma~\ref{lem:strongdvd}), which for a
  nowhere-vanishing sequence forces every Möbius quotient to be integral
  (Lemma~\ref{lem:moebiusfactor}) --- this replaces the periodicity computation above.
  Coprimality does use the valuation shape, but in the form ``$v_p(c_n) = E \cdot [m \dvd n]$ for
  some $m \ge 1$, $E$'' (\lname{cSeq\_factorization\_shape}), which is proved over an arbitrary
  UFD (Theorem~\ref{thm:valshape}) and fed to the general coprimality lemma
  \lname{moebiusFactorR\_isRelPrime}.}
\end{proof}

The next lemma gives us a simple irreducibility test for the $f_n$.

\begin{lemma}
  \label{lem:1.2}
  \uses{def:cseq, def:iterates}
  \lean{QuadraticIterates.irreducible_iteratedPoly_of_not_isSquare_cSeq,
    QuadraticIterates.even_reducible_factorization,
    QuadraticIterates.isUnit_or_associated_of_dvd_comp,
    QuadraticIterates.isUnit_or_associated_of_dvd_comp_of_associated,
    QuadraticIterates.map_iteratedPoly_comp_neg_X, QuadraticIterates.map_iteratedPoly_succ}
  \leanok
  If none of $c_1, c_2, \dots, c_n$ is a square in $\QQ$, then $f_n$ is irreducible in $\QQ[X]$.
\end{lemma}
\begin{proof}
  \leanok
  \uses{lem:evenpoly, lem:involution}
  Assume $f_n$ reducible, and let $m := \min\{k \mid f_k \text{ is reducible}\}$. Then
  $1 \le m \le n$, $f_{m-1}$ is irreducible, and $f_m = f_{m-1}(X^2 + a)$ is reducible. Let
  $\varepsilon := 1$ if $m > 1$ and $\varepsilon := -1$ if $m = 1$. Now $f_m$ is an even
  polynomial without nontrivial even divisors (since $f_{m-1}(Y + a)$ would be reducible).
  Therefore, if $h$ is a nontrivial divisor of $f_m$, then so is $h(-X)$, and $h$ and $h(-X)$ are
  coprime (since their g.c.d.\ is even). Thus, $f_m$ being reducible, there exists a polynomial
  $g$ with $\varepsilon \cdot f_m = g(X) \cdot g(-X)$. In particular,
  $c_m = \varepsilon f_m(0) = g(0)^2$ is a square in $\QQ$.
  \formalnote{The sign $\varepsilon$ is carried explicitly: what is produced is
  $(-1)^{\deg g} f_m = g \cdot g(-X)$ with $2 \deg g = \deg f_m$, and $(-1)^{\deg g}$ is exactly
  the sign relating $c_m$ to $f_m(0)$ (\lname{intCast\_cSeq\_eq\_ite\_mul\_eval\_zero}).
  ``No nontrivial even divisor'' is proved in the sharper form needed here: any divisor of
  $F \circ (X^2 + c)$ that is even \emph{up to associates} is a unit or an associate of the whole
  polynomial (\lname{isUnit\_or\_associated\_of\_dvd\_comp\_of\_associated}); the pairing of the
  irreducible factors under $X \mapsto -X$ is a fixed-point-free involution argument on the
  multiset of normalized factors (Lemma~\ref{lem:involution}).}
\end{proof}

\begin{corollary}
  \label{cor:1.3}
  \uses{lem:1.2}
  \lean{QuadraticIterates.irreducible_iteratedPoly_of_pos,
    QuadraticIterates.irreducible_iteratedPoly, QuadraticIterates.not_isSquare_cSeq,
    QuadraticIterates.not_isSquare_neg_of_irreducible}
  \leanok
  All $f_n$ are irreducible over $\QQ$.
\end{corollary}
\begin{proof}
  \leanok
  \uses{lem:1.1}
  The proof of Lemma~\ref{lem:1.1}~a) shows $|c_n| \ge |a|$ for all $n \ge 1$. Hence we have if
  $a > 0$: $c_n^2 < c_n^2 + a = c_{n+1} < (c_n + 1)^2$, so $c_{n+1}$ is not a square; if
  $a < -1$: $(c_n - 1)^2 < c_n^2 + a = c_{n+1} \le c_n^2$, so $c_{n+1}$ is not a square. Since
  $c_1 = -a$ is not a square either, Lemma~\ref{lem:1.2} applies.
  \formalnote{The formal statement covers $n = 0$ as well ($f_0 = X$ is irreducible). The
  converse implication is also formalized: irreducibility of a single $f_n$ with $n \ge 1$ forces
  $-a$ to be a non-square (\lname{not\_isSquare\_neg\_of\_irreducible}), since otherwise
  $f_n = (f_{n-1} - r)(f_{n-1} + r)$. This is what allows Lemma~\ref{lem:1.6} to be stated with
  irreducibility as its hypothesis.}
\end{proof}

The well-known example $a = -2$, where $K_n = \QQ(\zeta_{2^{n+2}}) \cap \mathbb{R}$, shows that
it is possible to have $\Om_n$ as small as $C_{2^n}$. We will see later that we can also give
examples where the other extreme, $\Om_n \cong \Ct{n}$, occurs.

\begin{lemma}
  \label{lem:1.4}
  \uses{def:galois, def:wreath, facts:1.0}
  \lean{QuadraticIterates.embed_equiv, QuadraticIterates.odoni_embedding,
    QuadraticIterates.isPGroup_galoisGroup, QuadraticIterates.pow_two_pow_apply_root,
    QuadraticIterates.ncard_rootSet_iteratedPoly_le,
    QuadraticIterates.nonempty_mulEquiv_of_finrank_eq}
  \leanok
  $\Om_n$ can be embedded into $\Ct{n}$. In particular, we have for all $n$:
  \[ \Om_{n+1} \cong \Ct{n+1} \iff \Om_n \cong \Ct{n} \text{ and } [K_{n+1} : K_n] = 2^{2^n}. \]
\end{lemma}
\begin{proof}
  \leanok
  \uses{lem:sylowwreath}
  Using the result of \cite{odoni}, p.~103, Cor.~1.2, and the fact that Galois groups do not
  increase under specialization, we get $\Om_n \hookrightarrow \Ct{n}$. The rest is then clear
  because of $\#\Ct{n+1}/\#\Ct{n} = 2^{2^n}$ and $[K_{n+1} : K_n] \le 2^{2^n}$
  (Facts~\ref{facts:1.0}).
  \formalnote{The embedding is proved directly rather than quoted, and without specializing from
  a generic polynomial. Every element of $\Om_n$ has order a power of $2$, because an
  automorphism raised to the power $2^m$ fixes each root of $f_m$ (induction on $m$, using that
  the roots of $f_{m+1}$ are the square roots of the shifted roots of $f_m$); so $\Om_n$ is a
  $2$-group acting faithfully on the at most $2^n$ roots of $f_n$, hence embeds into a Sylow
  $2$-subgroup of $S_{2^n}$, which is $\Ct{n}$ (Lemma~\ref{lem:sylowwreath}). The degree
  bookkeeping in the second part is done with $\#\Om_n = [K_n:\QQ]$ and
  $[K_{n+1}:\QQ] = [K_{n+1}:K_n]\,[K_n:\QQ]$.}
\end{proof}

\begin{lemma}
  \label{lem:1.5}
  \uses{def:twoindep, def:galois, def:cseq}
  \lean{QuadraticIterates.kummer_extension_criterion,
    QuadraticIterates.isSquare_algebraMap_iff_exists_mul_prod,
    QuadraticIterates.isSquare_algebraMap_iff_exists_sq_eq,
    QuadraticIterates.finrank_adjoin_range_eq_two_pow,
    QuadraticIterates.finrank_adjoin_le_two_pow,
    QuadraticIterates.card_monoidHom_eq_two_pow,
    QuadraticIterates.isSquare_algebraMap_cSeq,
    QuadraticIterates.twoIndependent_snoc_iff,
    QuadraticIterates.finrank_eq_of_nonempty_mulEquiv}
  \leanok
  If $\Om_n \cong \Ct{n}$ and $c_1, c_2, \dots, c_n$ are $2$-independent, then for all
  $c \in \QQ^*$:
  \[ c \notin (K_n^*)^2 \iff c_1, c_2, \dots, c_n, c \text{ are $2$-independent.} \]
\end{lemma}
\begin{proof}
  \leanok
  \uses{lem:wreathab, lem:multiquadratic, lem:descent}
  One proves easily by induction (using the fact that for a transitive permutation group $G$ one
  has $(G[H])^{ab} \cong G^{ab} \times H^{ab}$, where $G[H]$ denotes the wreath product) that
  $(\Ct{n})^{ab}$, the largest abelian factor group of $\Ct{n}$, is isomorphic with $C_2^n$.
  Thus, the largest $2$-Kummer extension of $\QQ$ within $K_n$ has degree $2^n$, and the kernel
  of the natural map $\QQ^*/(\QQ^*)^2 \to K_n^*/(K_n^*)^2$ has $\FF_2$-dimension $n$. Now,
  denoting by $A$ the set of zeros of $f_{m-1}$, we have that
  $c_m = f_{m-1}(-a) = \prod_{\alpha \in A} (-a - \alpha) = \prod_{\alpha \in A} (-a + \alpha)$
  is a square in $K_m$ (note that $A = -A$ and that for $\alpha \in A$, $\alpha - a$ is a square
  in $K_m$), so all of $c_1, \dots, c_n$ are squares in $K_n$. Since $c_1, \dots, c_n$ are
  $2$-independent, the above kernel must be generated by (the residue classes of)
  $c_1, \dots, c_n$. This implies the asserted equivalence.
  \formalnote{The two halves of ``the kernel has dimension exactly $n$'' are separate lemmas. For
  the lower bound, $\QQ(\sqrt{c_1}, \dots, \sqrt{c_n})$ has degree $2^n$ because the $c_i$ are
  $2$-independent (\lname{finrank\_adjoin\_range\_eq\_two\_pow}, via
  Lemma~\ref{lem:multiquadratic}). For the upper bound, any subfield of $K_n$ generated over
  $\QQ$ by square roots of rationals is Galois with elementary abelian group, so its degree is
  bounded by $\#\Hom(\Om_n, C_2) = 2^n$ (Lemma~\ref{lem:wreathab}). Hence a square root of $c$
  lying in $K_n$ already lies in $\QQ(\sqrt{c_1}, \dots, \sqrt{c_n})$, and the multiquadratic
  descent lemma (Lemma~\ref{lem:descent}) turns that into the statement that $c$ times a
  subproduct of the $c_i$ is a square in $\QQ$ --- which is exactly the failure of
  $2$-independence of $c_1, \dots, c_n, c$.}
\end{proof}

% The paper's remark after Lemma 1.5 is unnumbered, and it stays unnumbered here: a `remark`
% environment would consume a number and shift Lemma 1.6 to 1.7.
\emph{Remark.} We need to take $c_1 = -a$ instead of the more natural $c_1 = f(0) = a$ in order
to have $c_1$ a square in $K_1$ (\lname{cSeq\_one}).

\begin{lemma}[cf.\ \cite{odoni}, Lemma 4.3]
  \label{lem:1.6}
  \uses{def:galois, def:cseq, facts:1.0}
  \lean{QuadraticIterates.degree_criterion, QuadraticIterates.relfinrank_succ_eq_pow,
    QuadraticIterates.rootShift, QuadraticIterates.coe_rootShift,
    QuadraticIterates.rootShift_ne_zero, QuadraticIterates.prod_radicand_eq_cSeq,
    QuadraticIterates.prod_aroots_sub_eq_cSeq,
    QuadraticIterates.isPretransitive_galoisGroup,
    QuadraticIterates.exists_ringEquiv_radicand_smul,
    QuadraticIterates.degree_criterion_zero,
    QuadraticIterates.finrank_splittingField_one}
  \leanok
  For all $n$: $[K_{n+1} : K_n] = 2^{2^n} \iff c_{n+1}$ is not a square in $K_n$.
\end{lemma}
\begin{proof}
  \leanok
  \uses{cor:1.3, lem:multiquadratic, lem:invariant, lem:relations}
  Let $\alpha_j$ $(j = 1, 2, \dots, 2^n)$ be the roots of $f_n$ in $K_n$. Then
  $K_n \subset K_{n+1}$ is a $2$-Kummer extension, and $[K_{n+1} : K_n] = 2^{2^n - \dim V}$,
  where $V$ is the $\FF_2$-vector space of all $(\varepsilon_1, \dots, \varepsilon_{2^n}) \in
  \FF_2^{2^n}$ such that $\prod_j (\alpha_j - a)^{\varepsilon_j}$ is a square in $K_n$. We show
  that $V \ne 0 \iff (1, \dots, 1) \in V$: $\Om_n$ operates on $V$ by permuting the components of
  the elements of $V$ according to its action on the $\alpha_j$, thus $V$ becomes an
  $\FF_2[\Om_n]$-module. Now, if $G$ is a $2$-group and $M \ne 0$ is an $\FF_2[G]$-module, then
  the submodule $M^G$ of $G$-invariant elements of $M$ is also nontrivial. (Induction on $\#G$:
  Let $G \cong C_2 = \{1, \sigma\}$ and $0 \ne y \in M$. Either $y = \sigma y$, then
  $0 \ne y \in M^G$, or $y \ne \sigma y$, then $0 \ne y + \sigma y \in M^G$. If $\#G > 2$, take a
  nontrivial normal subgroup $N$ of $G$. Then $M$ is an $\FF_2[N]$-module, too, so by induction
  $M^N \ne 0$. Now $M^N$ is an $\FF_2[G/N]$-module, so $M^G = (M^N)^{G/N} \ne 0$.) So, if
  $V \ne 0$, then $V^{\Om_n} \ne 0$. Since $\Om_n$ operates transitively on the $\alpha_j$
  ($f_n$ is irreducible by Cor.~\ref{cor:1.3}), the only nontrivial invariant element in $V$ is
  $(1, \dots, 1)$.

  From the above it is clear that the left hand side is equivalent to
  $(1, \dots, 1) \notin V$, whereas $(1, \dots, 1) \in V$ means that
  $\prod_j (\alpha_j - a) = c_{n+1}$ is a square in $K_n$.
  \formalnote{The hypothesis used is that $f_n$ is irreducible, which by
  Corollary~\ref{cor:1.3} holds under the standing assumption on $a$; the paper instead argues
  under the maximality hypothesis $[K_n : \QQ] = 2^{2^n - 1}$, which is strictly stronger. The
  degenerate case $n = 0$ is included ($K_1 = \QQ(\sqrt{-a})$ has degree $2$ over $\QQ$ exactly
  when $-a$ is not a square). $V$ is \texttt{rootRelations} of the family
  $\beta \mapsto \beta - a$ of shifted roots (\texttt{rootShift}); that this family consists of
  nonzero elements is where $c_{n+1} \neq 0$ enters. The identity
  $\prod_j (\alpha_j - a) = c_{n+1}$ is proved by evaluating the monic split $f_n$ at $a$
  (\lname{prod\_aroots\_sub\_eq\_cSeq}).}
\end{proof}

Summing up, we can now prove the following

\begin{theorem}[Section 1]
  \label{thm:section1}
  \uses{def:galois, def:wreath, def:twoindep, def:cseq, def:bseq}
  \lean{QuadraticIterates.section1_equiv, QuadraticIterates.section1_a_iff_b,
    QuadraticIterates.section1_b_iff_c, QuadraticIterates.section1_squarefree,
    QuadraticIterates.sqClass_cSeq_eq_sum_divisors,
    QuadraticIterates.sqClass_bSeq_eq_sum_divisorsAntidiagonal}
  \leanok
  \begin{enumerate}
    \item For all $n$, the following statements are equivalent:
      \begin{enumerate}
        \item[a)] $\Om_n \cong \Ct{n}$;
        \item[b)] $c_1, c_2, \dots, c_n$ are $2$-independent;
        \item[c)] $b_1, b_2, \dots, b_n$ are $2$-independent.
      \end{enumerate}
    \item If none of $|b_2|, |b_3|, \dots, |b_n|$ is a square in $\QQ$, then $\Om_n \cong \Ct{n}$.
  \end{enumerate}
\end{theorem}
\begin{proof}
  \leanok
  \uses{lem:1.1, lem:1.4, lem:1.5, lem:1.6, cor:1.3, lem:coprimesquare}
  1) By induction on $n$. $n = 0$ is trivial. So let $n \ge 1$ and suppose 1) holds for $n - 1$.

  ``a) $\Leftrightarrow$ b)'': $\Om_n \cong \Ct{n}$ $\iff$ (by \ref{lem:1.4})
  $\Om_{n-1} \cong \Ct{n-1}$ and $[K_n : K_{n-1}] = 2^{2^{n-1}}$ $\iff$ (by induction hypothesis
  and \ref{lem:1.6}) $c_1, \dots, c_{n-1}$ are $2$-independent and $c_n \notin (K_{n-1}^*)^2$
  $\iff$ (by \ref{lem:1.5} and induction hypothesis) $c_1, \dots, c_n$ are $2$-independent.

  ``b) $\Leftrightarrow$ c)'': This is clear because of $c_m = \prod_{d \dvd m} b_d$.

  2) Since by Lemma~\ref{lem:1.1} $b_1, b_2, \dots, b_n$ are pairwise coprime, 1)~c) is
  equivalent to: none of the $b_m$ is a square, and at most one of the $-b_m$ is a square. Since
  $b_1 = -a$ is not a square, condition 2) implies 1)~c).
  \formalnote{In 1), ``b) $\Leftrightarrow$ c)'' is run in $\QQ^*/(\QQ^*)^2$: the classes of the
  $c_m$ and of the $b_m$ span the same $\FF_2$-subspace (each is a $\ZZ$-combination of the
  others, by Möbius inversion in both directions), and $2$-independence of a family of $n$
  vectors is equivalent to that span having dimension $n$. In 2), the passage from a square
  subproduct to a single square factor is the pairwise coprimality of the $b_m$: a square product
  of pairwise coprime integers has all $|b_m|$ square (Lemma~\ref{lem:coprimesquare}); the index
  $m = 1$ is excluded separately using $b_1 = -a$.}
\end{proof}

\chapter{Iteration sequences associated to even integer polynomials}

% The counter is set so that the two lemmas below carry the numbers 2.1, 2.2 of the paper.
\setcounter{theorem}{-1}

In this section, let $g \in \ZZ[X^2]$ be an even polynomial with integer coefficients. Put
$\gamma_1 := \pm g(0)$, and for $n \ge 1$, $\gamma_{n+1} := g(\gamma_n)$. Assume that all
$\gamma_n \ne 0$ and set $\beta_n := \prod_{d \dvd n} \gamma_d^{\mu(n/d)}$ for $n \ge 1$.

\begin{definition}[The sequences $\gamma$ and $\beta$]
  \label{def:gammaseq}
  \lean{QuadraticIterates.EvenPoly, QuadraticIterates.evenPoly_iff_comp_neg_X,
    QuadraticIterates.EvenPoly.eval_neg, QuadraticIterates.EvenPoly.eval_congr,
    QuadraticIterates.EvenPoly.dvd_eval_sub, QuadraticIterates.EvenPoly.map,
    QuadraticIterates.EvenPoly.iterate_eval_neg,
    QuadraticIterates.gammaSeq, QuadraticIterates.gammaSeq_zero,
    QuadraticIterates.gammaSeq_one, QuadraticIterates.gammaSeq_succ,
    QuadraticIterates.gammaSeq_add, QuadraticIterates.map_gammaSeq,
    QuadraticIterates.intCast_gammaSeq, QuadraticIterates.betaSeq,
    QuadraticIterates.betaSeq_eq_moebiusFactorR, QuadraticIterates.intCast_betaSeq}
  \leanok
  $\gamma_1 = \varepsilon\,g(0)$ with $\varepsilon = \pm 1$, $\gamma_{n+1} = g(\gamma_n)$ for
  $n \ge 1$, and $\beta_n = \prod_{d \dvd n} \gamma_d^{\mu(n/d)}$.
  \formalnote{Being even is defined as membership in $R[X^2]$, i.e.\ $g = h(X^2)$ for some $h$;
  over a domain of characteristic $\ne 2$ this is equivalent to $g(-X) = g$
  (\lname{evenPoly\_iff\_comp\_neg\_X}), and both forms are used. As with $c$, the sequence
  $\gamma$ is total with junk value $\gamma_0 = 0$, and $\beta_n$ is an element of the
  coefficient ring, defined as the preimage of the Möbius product in the fraction field. The
  sequences are set up over an arbitrary commutative semiring, and each result below is stated
  over the weakest structure that carries it: a commutative ring for the congruences, a GCD
  domain for strong divisibility, a UFD for the valuation shape, and $\ZZ$ for Lemmas~2.1
  and~2.2. Since the whole construction commutes with ring homomorphisms
  (\lname{map\_gammaSeq}), congruences mod $m$ are proved by computing in $\ZZ/m\ZZ$, and $c$ is
  literally the case $g = X^2 + a$, $\varepsilon = -1$.}
\end{definition}

\begin{lemma}
  \label{lem:2.1}
  \uses{def:gammaseq}
  \lean{QuadraticIterates.not_isSquare_betaSeq, QuadraticIterates.gammaSeq_mul_eq_two_mul,
    QuadraticIterates.prod_gammaSeq_mul_eq, QuadraticIterates.prod_gammaSeq_mul_eq_neg_prod,
    QuadraticIterates.isUnit_prod_gammaSeq_mul, QuadraticIterates.gammaSeq_gcd,
    QuadraticIterates.gammaSeq_associated_gcd, QuadraticIterates.gammaSeq_associated_one,
    QuadraticIterates.gammaSeq_eq_gammaSeq_one}
  \leanok
  Suppose that for all $n \ge 1$ there is an $m$ such that $m \dvd \gamma_n + \gamma_{2n}$, $m$
  is prime to $\gamma_n$, and $-1$ is not a square mod $m$. Then for all $n \ge 2$, $\beta_n$ is
  not a square in $\QQ$.
\end{lemma}
\begin{proof}
  \leanok
  \uses{def:bseq, lem:moebiussign, lem:strongdvd, lem:zmodsquare}
  Let $n > 1$ have prime decomposition $n = p_1^{e_1} \cdots p_r^{e_r}$, where $r \ge 1$, the
  $p_j$ are distinct primes, and all $e_j \ge 1$. Set $n' := p_1 \cdots p_r > 1$ and
  $k := n/n'$. Put $\alpha := \gamma_{2k}$. Let $m$ be such that
  $m \dvd \gamma_k + \gamma_{2k}$, $m$ is prime to $\gamma_k$ and hence to $\alpha$, and $-1$ is
  not a square mod $m$. Then $\gamma_k \equiv -\alpha \bmod m$, and
  $\gamma_{3k} = g_k(\gamma_{2k}) = g_k(-\gamma_k) = g_k(\gamma_k) = \gamma_{2k} = \alpha
  \bmod m$ --- writing $g_k$ for the $k$-fold iterate of $z \mapsto g(z)$, which is an even
  function --- and by induction $\gamma_{tk} \equiv \alpha \bmod m$ for all $t \ge 2$. Therefore
  \[ \beta_n = \prod_{d \dvd n} \gamma_d^{\mu(n/d)}
             = \prod_{t \dvd n'} \gamma_{kt}^{\mu(n'/t)}
             \equiv (-1)^{\mu(n')} \prod_{t \dvd n'} \alpha^{\mu(n'/t)}
             = -1 \bmod m \]
  (for $\mu(n') = \pm 1$ and $\sum_{t \dvd n'} \mu(n'/t) = 0$). Since $-1$ is not a square mod
  $m$, $\beta_n$ cannot be a square in $\QQ$.
  \formalnote{$n'$ is $\rad(n)$, and the Möbius product is split along the partition of the
  divisors of $n'$ into those with $\mu = +1$ and those with $\mu = -1$
  (Lemma~\ref{lem:moebiussign}), which have equal cardinality; $\beta_n$ is then the quotient
  $P/Q$ of the two subproducts, and the computation above becomes $P \equiv -Q \bmod m$ with $Q$
  a unit mod $m$. That a quotient of two such integers cannot be a square when $-1$ is not a
  square mod $m$ is \lname{ZMod.isSquare\_neg\_one\_of\_isSquare\_div}
  (Lemma~\ref{lem:zmodsquare}). The sign $(-1)^{\mu(n')}$ comes from the single index $t = 1$,
  where $\gamma_k \equiv -\alpha$ rather than $\alpha$; this bookkeeping is
  \lname{prod\_gammaSeq\_mul\_eq}.}
\end{proof}

We will now give a handier criterion for the assumptions of Lemma~\ref{lem:2.1} to hold.

\begin{lemma}
  \label{lem:2.2}
  \uses{def:gammaseq, lem:2.1}
  \lean{QuadraticIterates.not_isSquare_betaSeq_of_pos,
    QuadraticIterates.not_isSquare_betaSeq_of_pos_of_not_isSquare_neg_one,
    QuadraticIterates.not_isSquare_betaSeq_of_pos_of_eval_one_emod_four_eq_two,
    QuadraticIterates.not_isSquare_betaSeq_of_pos_of_eval_one_emod_four_eq_three,
    QuadraticIterates.gammaSeq_one_eq_one_of_pos,
    QuadraticIterates.gammaSeq_add_succ_dvd, QuadraticIterates.isCoprime_gammaSeq_add_succ,
    QuadraticIterates.gammaSeq_add_succ_dvd_sub,
    QuadraticIterates.gammaSeq_add_succ_zmod_four_eq_three,
    QuadraticIterates.gammaSeq_add_succ_zmod_eight_eq_six,
    QuadraticIterates.gammaSeq_eq_ite_even, QuadraticIterates.gammaSeq_add_succ_eq_three,
    QuadraticIterates.gammaSeq_eq_eval_one,
    QuadraticIterates.gammaSeq_add_succ_eq_two_mul_eval_one}
  \leanok
  If all the $\gamma_n$ are positive, and
  \begin{enumerate}
    \item[a)] $g(0) = 1$ and $g(1) \equiv 2 \bmod 4$, or
    \item[b)] $g(0) = \pm 1$ and $g(1) \equiv 3 \bmod 4$,
  \end{enumerate}
  then $\beta_n$ is not a square in $\QQ$ for $n \ge 2$.
\end{lemma}
\begin{proof}
  \leanok
  \uses{lem:zmodsquare}
  Notice first that since $\gamma_n + \gamma_{n+1} \dvd \gamma_n + \gamma_{2n}$ (same argument as
  in the proof of Lemma~\ref{lem:2.1}, applied $(n-1)$ times), and $\gamma_n$ and
  $\gamma_{n+1}$ are coprime (since $\gamma_{n+1} = g(\gamma_n) \equiv g(0) = \pm 1
  \bmod \gamma_n$), we can take $m := \gamma_n + \gamma_{n+1}$ in Lemma~\ref{lem:2.1}, if $-1$ is
  not a square mod $m$. This holds for $m \equiv 3 \bmod 4$ and for $m \equiv 6 \bmod 8$.

  a) $\gamma_n \equiv 1$ or $2 \bmod 4$, according to whether $n$ is odd or even, so for all $n$,
  $\gamma_n + \gamma_{n+1} \equiv 3 \bmod 4$.

  b) $\gamma_1 + \gamma_2 \equiv 0 \bmod 4$, and for $n \ge 2$,
  $\gamma_n + \gamma_{n+1} \equiv 6 \bmod 8$, since all $\gamma_n$ $(n \ge 2)$ are
  $\equiv 3 \bmod 4$ and equal mod $8$.
  \formalnote{Positivity of the $\gamma_n$ is what makes $m := \gamma_n + \gamma_{n+1}$ a
  \emph{positive} integer, so that ``$-1$ is not a square mod $m$'' is a statement about
  $\ZZ/m\ZZ$. The two congruence computations are done in $\ZZ/4\ZZ$ resp.\ $\ZZ/8\ZZ$, over an
  arbitrary ring: for a) the only input is $4 = 0$, $g(0) = 1$, $g(1) = 2$, which forces the
  sequence to alternate between $1$ and $2$; for b) the only input is that the relevant residues
  square to $1$, which forces $\gamma_n = g(1)$ for all $n \ge 2$. The case $n = 1$ of b) is the
  separate observation that $4 \dvd \gamma_1 + \gamma_2$ and that $-1$ is not a square mod $4$.}
\end{proof}

In the next section, we will apply this result to the Galois groups of Chapter~1.

\chapter{\texorpdfstring{Some values of $a$ with $\Om_n \cong \Ct{n}$ for all $n$}{Some values of
  a with the full wreath power for all n}}

We combine the results of Chapters~1 and~2 to get the following.

\begin{definition}[The rescaled polynomial]
  \label{def:normpoly}
  \uses{def:cseq, def:gammaseq}
  \lean{QuadraticIterates.normPoly, QuadraticIterates.eval_normPoly,
    QuadraticIterates.evenPoly_normPoly, QuadraticIterates.evenPoly_X_sq_add_C,
    QuadraticIterates.evenPoly_C_mul_X_sq_add_C,
    QuadraticIterates.gammaSeq_normPoly_one, QuadraticIterates.abs_cSeq_eq_gammaSeq_mul_abs,
    QuadraticIterates.gammaSeq_normPoly_pos, QuadraticIterates.abs_bSeq_eq_betaSeq,
    Int.sign_sq_of_ne_zero}
  \leanok
  $g := |a|\,X^2 + \sgn(a)$, an even polynomial with $\gamma_1 = \sgn(a)\,g(0) = 1$. Its
  iteration sequence satisfies $\gamma_n \cdot |a| = |c_n|$ for all $n$ (both sides vanish at
  $n = 0$); in particular $\gamma_n > 0$ for $n \ge 1$, and $\beta_n = |b_n|$ for $n \ge 2$.
  \formalnote{Substituting $x \mapsto |a|\,x$ turns the recursion $c_{n+1} = c_n^2 + a$ into the
  $\gamma$-recursion of $g$, which is why the rescaled polynomial is introduced: $X^2 + a$ itself
  does not have a positive iteration sequence when $a < 0$, and positivity is what
  Lemma~\ref{lem:2.2} needs. The identity $\beta_n = |b_n|$ follows from
  $\gamma_n = |c_n|/|a|$ because the exponents $\mu(n/d)$ sum to zero for $n \ge 2$, so the
  factors $|a|$ cancel.}
\end{definition}

\begin{theorem}[Section 3]
  \label{thm:section3}
  \uses{def:normpoly, def:galois, def:wreath}
  \lean{QuadraticIterates.section3_main, QuadraticIterates.not_isSquare_abs_bSeq,
    QuadraticIterates.not_isSquare_neg_of_cases}
  \leanok
  If $a \in \ZZ$ has one of the following properties:
  \begin{enumerate}
    \item[a)] $a > 0$ and $a \equiv 1 \bmod 4$;
    \item[b)] $a > 0$ and $a \equiv 2 \bmod 4$;
    \item[c)] $a < 0$, $a \equiv 0 \bmod 4$, and $-a$ is not a square in $\ZZ$;
  \end{enumerate}
  then, denoting by $f_n$ the $n$-th iterate of $f := X^2 + a$, we have
  \[ \Gal(f_n/\QQ) \cong \Ct{n} \quad\text{for all } n \ge 1. \]
\end{theorem}
\begin{proof}
  \leanok
  \uses{thm:section1, lem:2.2, lem:1.1}
  In all cases, $-a$ is not a square in $\ZZ$. Let $g := |a|\,X^2 + \sgn(a)$. Then we have in
  case a) $g(0) = 1$ and $g(1) \equiv 2 \bmod 4$, and in cases b) and c) $g(1) \equiv 3 \bmod 4$.
  We put $\gamma_1 := +1$ and have $\gamma_n \cdot |a| = c_n$ for all $n \ge 2$, therefore
  $\gamma_n > 0$ for all $n \ge 1$, cf.\ Lemma~\ref{lem:1.1}~a). Now, according to
  Lemma~\ref{lem:2.2}, for $n \ge 2$, $|b_n| = \beta_n$ is not a square in $\QQ$. We can thus
  apply part 2) of Theorem~\ref{thm:section1}, which gives us $\Gal(f_n/\QQ) \cong \Ct{n}$ for
  all $n$.
  \formalnote{The three cases are stated in Lean exactly as here, as
  \texttt{(0 < a $\wedge$ a \% 4 = 1) $\vee$ (0 < a $\wedge$ a \% 4 = 2) $\vee$
  (a < 0 $\wedge$ a \% 4 = 0 $\wedge$ $\neg$IsSquare (-a))}. In case b) the value $g(1) =
  |a| + 1 \equiv 3 \bmod 4$, and in case c) $g(1) = |a| - 1 \equiv 3 \bmod 4$; case a) satisfies
  hypothesis a) of Lemma~\ref{lem:2.2} and the other two hypothesis b). The formal proof follows
  the paper's sentences: \lname{not\_isSquare\_neg\_of\_cases} (in all cases $-a$ is not a
  square), \lname{not\_isSquare\_abs\_bSeq} (Lemma~\ref{lem:2.2} for $g$, so no $|b_n|$ with
  $n \ge 2$ is a square), and part 2) of Theorem~\ref{thm:section1}.}
\end{proof}

\section*{Concluding remarks}

Odoni \cite{odoni} investigates the case $f = X^2 + 1$, i.e.\ $a = 1$. His result
(\cite{odoni}, p.~111, L.~4.4~(iv)) is equivalent in this case to Theorem~\ref{thm:section1}. He
then gives an algorithm deciding whether $\Om_n \cong \Ct{n}$ for any given $n$, which has been
carried out by Cremona \cite{cremona} up to $n = 5 \cdot 10^7$, leading him to conjecture that
$|b_n|$ for $n > 1$ is never a square, i.e.\ $\Om_n \cong \Ct{n}$ for all $n$. This conjecture is
now a special case of Theorem~\ref{thm:section3}, part a).

The more general problem of determining all $a \in \ZZ$ such that $\Om_n \cong \Ct{n}$ for all
$n$ remains open, as does the still more general problem of determining the groups $\Om_n$ when
$a \in \ZZ$ is given. Besides the cases dealt with in this paper and the (exceptional) case
$a = -2$, where $\Om_n \cong C_{2^n}$, not much seems to be known in this direction.

However, it is clear that there are infinitely many $a$ such that $\#\Om_2 = 4$ (and therefore
not for all $n$, $\Om_n \cong \Ct{n}$), since this is equivalent to $b_1$ and $b_2$ being
$2$-dependent, meaning $(-b_1)$ being a square, i.e.\ $b_2 = b^2$ is a square (since $b_1 = -a$
and $b_1 b_2 = a^2 + a$ cannot be squares), i.e.\ $a = -b^2 - 1$ with some $b > 0$. Moreover,
since $b_1 b_3 = c_3 = c_2^2 + a$ is never a square and (as is easily seen) $b_2$ and $b_3$ are
both squares only if $a = -2$, we have $[K_2 : \QQ] = 4$ and $[K_3 : K_2] < 16$ only for
$a = -2$; so in particular, $-2$ is the only value such that $\Om_2 \cong C_4$ and
$\Om_3 \cong C_8$ (or such that for all $n$, $\Om_n \cong C_{2^n}$).

A final remark concerning the sequence $(1, 2, 5, 26, 677, \dots)$ discussed in \cite{cremona}
and \cite{odoni}, which comes up in the case $a = 1$: numerical calculations reveal $603$ primes
up to $11\,684\,149$ which divide one of the $b_n$, but none of them divides ``its'' $b_n$ more
than once. Thus, one may be led to conjecture that the $b_n$, and therefore the $c_n$ too, are
squarefree. To prove this (if true) seems to be quite hard in any case.

\emph{Note added in proof.} There are plausibility arguments indicating that for a given prime
$p$, the ``probability'' that there is some $n$ such that $p^2 \dvd b_n$ should be approximately
$\sqrt{\pi/(8p^3)}$. Considering the above data, the ``probability'' for the conjecture to be
true should then be at least $\prod_{p > 11\,684\,149} \bigl(1 - \sqrt{\pi/(8p^3)}\bigr)$, which
is very near $1$.

\emph{None of the material in these concluding remarks is formalized}: the statements about
$a = -2$, about $\#\Om_2 = 4$ and the squarefreeness heuristics are outside the scope of the
formalization, which covers the numbered results of the three sections.

\chapter{Auxiliary theory}

The paper carries out several arguments by hand that are really instances of general theory:
degrees of multiquadratic extensions, invariants of a $2$-group acting on an $\FF_2$-vector
space, abelian quotients of iterated wreath products, and Möbius quotients of a strong
divisibility sequence. In the formalization these are developed separately, in the generality
they naturally have, under \texttt{QuadraticIterates/Mathlib/}; they are candidates for
upstreaming into Mathlib. This chapter records them, since the proofs of Chapters~1--3 refer to
them.

\section{Multiquadratic extensions}

\begin{lemma}[Degree of a multiquadratic extension]
  \label{lem:multiquadratic}
  \lean{multiquadraticRelations, mem_multiquadraticRelations, multiquadratic_degree,
    multiquadratic_degree_family, multiquadratic_degree_insert_of_maximal,
    multiquadratic_degree_insert_family, finrank_adjoin_sq_le, finrank_adjoin_finset_sq_le,
    finrank_adjoin_sq_eq, multiquadraticRelations_finrank_le,
    multiquadraticRelations_insert_finrank}
  \leanok
  Let $F$ be a field of characteristic $\ne 2$ and $s$ a finite set of square roots of elements
  $c_y \in F^*$ inside an extension. Then
  $[F(s) : F] = 2^{\#s - \dim_{\FF_2} R(s)}$, where $R(s) \subseteq \FF_2^{s}$ is the space of
  \emph{relations}: the $\varepsilon$ with $\prod_{y : \varepsilon_y = 1} c_y$ a square in $F$.
  In particular $[F(s) : F] \le 2^{\#s}$, with equality iff there are no relations; and adjoining
  one further square root either doubles the degree or leaves it unchanged.
\end{lemma}
\begin{proof}
  \leanok
  Induction on $\#s$, the step being the descent Lemma~\ref{lem:descent}: the degree doubles
  exactly when the new radicand is not a square in the field generated so far.
\end{proof}

\begin{lemma}[Square descent]
  \label{lem:descent}
  \lean{square_descent, square_descent_step, isSquare_algebraMap_iff,
    not_isSquare_algebraMap_of_sqrt_notMem, isSquare_algebraMap_bot_iff, mem_adjoin_simple_sq,
    sqClass_mul, sqClass_prod, sqClass_prod_zpow, sqClass_eq_zero_iff,
    isSquare_prod_iff_sum_sqClass_eq_zero}
  \leanok
  An element of $F^*$ is a square in a multiquadratic extension $F(\sqrt{c_y} : y \in s)$ if and
  only if it becomes a square in $F$ after multiplication by a subproduct of the $c_y$.
\end{lemma}
\begin{proof}
  \leanok
  Adjoining one square root at a time; in each step, an element of the smaller field whose square
  root lies in the quadratic extension differs from a square by the radicand.
\end{proof}

\begin{lemma}[Relations among the shifted roots]
  \label{lem:relations}
  \lean{rootRelations, mem_rootRelations, rootRelations_finrank_le,
    rootRelations_finrank_reindex, multiquadraticRelations_finrank_eq_rootRelations,
    all_ones_mem_rootRelations, extendByZeroLM, extendByZeroLM_injective,
    multiquadraticRelations_ycoord, multiquadraticRelations_ker_finrank,
    multiquadraticRelations_finite}
  \leanok
  For an injective family $(r_i)_{i \in \iota}$ of nonzero elements of $F$, the relation space
  $R \subseteq \FF_2^{\iota}$ indexed by the family agrees with the relation space of its image
  set, and $\dim R \le \#\iota$. The all-ones vector lies in $R$ if and only if $\prod_i r_i$ is
  a square in $F$.
\end{lemma}
\begin{proof}
  \leanok
  Reindexing along the bijection $\iota \to \mathrm{range}$, plus the definition of the relation
  space.
\end{proof}

\begin{lemma}[Invariants of a $2$-group]
  \label{lem:invariant}
  \lean{fixed_points_nontrivial, invariant_submodule_all_ones, rootRelations_invariant,
    rootRelations_all_ones}
  \leanok
  If a finite $2$-group $G$ acts on a nontrivial $\FF_2$-vector space $M$, then $M^G \ne 0$. If
  moreover $M \subseteq \FF_2^{\iota}$ is a submodule stable under a \emph{transitive} action of
  $G$ on $\iota$ permuting coordinates, then $M \ne 0$ implies that the all-ones vector lies in
  $M$.
\end{lemma}
\begin{proof}
  \leanok
  For the first statement, $\#M$ is a power of $2$, and $\#M \equiv \#M^G \bmod 2$ by Mathlib's
  \lname{IsPGroup.card\_modEq\_card\_fixedPoints}; hence $\#M^G > 1$.
  For the second, a nonzero invariant vector has constant coordinates along the orbits of $G$ on
  $\iota$, and there is only one orbit.
\end{proof}

\section{Wreath products}

\begin{lemma}[Sylow $2$-subgroups of symmetric groups]
  \label{lem:sylowwreath}
  \lean{Sylow.mulEquivIteratedWreathProduct, IteratedWreathProduct.card,
    MonoidHom.nonempty_mulEquiv_iff_card_eq}
  \leanok
  A Sylow $2$-subgroup of the symmetric group on $2^n$ letters is isomorphic to $\Ct{n}$, which
  has order $2^{2^n - 1}$. Consequently, an injective homomorphism between finite groups of equal
  order is an isomorphism, which is how ``$\Om_n \cong \Ct{n}$'' is obtained from an embedding
  plus a degree count.
\end{lemma}
\begin{proof}
  \leanok
  Both statements are in Mathlib (\texttt{Sylow.mulEquivIteratedWreathProduct}); the last one is
  elementary.
\end{proof}

\begin{lemma}[The abelian quotient of $\Ct{n}$]
  \label{lem:wreathab}
  \lean{wreath_max_elem_ab, RegularWreathProduct.card_hom, elem_ab_card_hom,
    isGalois_adjoin_of_sq_eq_algebraMap, algEquiv_adjoin_sq_eq_one,
    apply_eq_or_eq_neg_of_sq_eq_algebraMap}
  \leanok
  $\#\Hom(\Ct{n}, C_2) = 2^n$; equivalently, the largest elementary abelian $2$-quotient of
  $\Ct{n}$ is $C_2^n$. A field extension generated by square roots of base-field elements is
  Galois with elementary abelian group of exponent $2$, so its degree equals the number of
  homomorphisms of its Galois group to $C_2$.
\end{lemma}
\begin{proof}
  \leanok
  For a regular wreath product $A \wr B$ the homomorphisms to an abelian group of exponent $2$
  factor through the coinvariants of the base, giving
  $\#\Hom(A \wr B, C_2) = \#\Hom(A, C_2) \cdot \#\Hom(B, C_2)$; induction on $n$ then gives
  $2^n$. The statement about elementary abelian groups is the identity $\#G = \#\Hom(G, C_2)$
  for such $G$.
\end{proof}

\section{Möbius quotients of a strong divisibility sequence}

\begin{lemma}[Strong divisibility from a translation congruence]
  \label{lem:strongdvd}
  \lean{associated_gcd_of_dvd_sub, gcd_eq_normalize_of_dvd_sub, Int.gcd_eq_natAbs_of_dvd_sub}
  \leanok
  Let $(c_n)$ be a sequence in a GCD domain with $c_0 = 0$ and $c_m \dvd c_{m+j} - c_j$ for all
  $m, j$. Then $\gcd(c_m, c_n)$ is associated to $c_{\gcd(m,n)}$.
\end{lemma}
\begin{proof}
  \leanok
  The Euclidean algorithm on the indices: the hypothesis gives
  $\gcd(c_{m+j}, c_j) = \gcd(c_m, c_j)$, which is the step of the subtractive algorithm; the
  junk value $c_0 = 0$ makes the base case $\gcd(c_m, c_0) = c_m$ come out right.
\end{proof}

\begin{lemma}[Integrality of the Möbius quotients]
  \label{lem:moebiusfactor}
  \lean{moebiusFactorK, moebiusFactorK_eq_div, moebiusFactorR, numProd, denProd,
    denProd_dvd_numProd, moebiusFactorK_isInteger, algebraMap_moebiusFactorR,
    prod_moebiusFactorR, moebiusFactorR_ne_zero, moebiusFactorR_mul_denProd,
    factorization_gcd_min, isInteger_div_iff_dvd}
  \leanok
  Let $R$ be a UFD and $(c_n)$ a sequence in $R \setminus \{0\}$ (on positive indices) which is a
  strong divisibility sequence. Then each Möbius quotient
  $b_n := \prod_{d \dvd n} c_d^{\mu(n/d)}$ lies in $R$ and is nonzero, and Möbius inversion
  holds: $c_n = \prod_{d \dvd n} b_d$.
\end{lemma}
\begin{proof}
  \leanok
  Write the quotient as $N/D$ with $N$ the product of the $c_d$ with $\mu(n/d) = 1$ and $D$ the
  product of those with $\mu(n/d) = -1$. At each prime $p$, strong divisibility bounds
  $v_p(c_{\gcd(d,d')})$ below by $\min(v_p(c_d), v_p(c_{d'}))$, and a counting argument on the
  divisor lattice shows $v_p(D) \le v_p(N)$; hence $D \dvd N$.
\end{proof}

\begin{theorem}[Constant valuation shape]
  \label{thm:valshape}
  \lean{factorization_periodic_shape, factorization_moebiusFactorR_shape,
    moebiusFactorR_isRelPrime, factorization_eq_of_dvd_sub, one_le_factorization_iff_dvd,
    factorization_eq_zero_iff_not_dvd, pow_dvd_iff_le_factorization,
    QuadraticIterates.factorization_gammaSeq_shape,
    QuadraticIterates.factorization_gammaSeq_of_dvd_eval_zero,
    QuadraticIterates.gammaSeq_period, QuadraticIterates.sq_dvd_gammaSeq_succ_sub,
    QuadraticIterates.pow_succ_dvd_gammaSeq_succ_sub}
  \leanok
  For the iteration sequence of an even polynomial with $\varepsilon^2 = 1$ over a UFD, and each
  normalized prime $p$, there are $m \ge 1$ and $E$ with $v_p(\gamma_n) = E$ for $m \dvd n$ and
  $v_p(\gamma_n) = 0$ otherwise. Consequently the valuation of each Möbius quotient $\beta_n$ is
  supported on the single index $n = m$, so the $\beta_n$ are pairwise coprime.
\end{theorem}
\begin{proof}
  \leanok
  Two cases. If $p \dvd g(0)$, the valuation is constantly $v_p(g(0))$, since
  $\gamma_n^2 \dvd \gamma_{n+1} - g(0)$ forces $v_p(\gamma_{n+1}) = v_p(g(0))$ by induction. If
  not, let $m$ be minimal with $p \dvd \gamma_m$ and $E := v_p(\gamma_m)$; then
  $p^{E+1} \dvd \gamma_{m+1} - g(0)$, whence $p \nmid \gamma_{m+1}$, and the sequence is periodic
  mod $p^{E+1}$ with period $m$ from index $2$ on --- this is the periodicity argument of the
  proof of Lemma~\ref{lem:1.1}~b), isolated and stated for an arbitrary sequence with these
  properties (\lname{factorization\_periodic\_shape}).
\end{proof}

\begin{lemma}[Möbius sums]
  \label{lem:moebiussign}
  \lean{moebius_sign_partition, prod_pow_moebius_eq_div, moebius_antidiag_sum_zero,
    beta_radical, moebius_restricted_sum, indicator_moebius_nonneg, moebius_transform_nonneg}
  \leanok
  For squarefree $n' > 1$, the divisors of $n'$ split into two sets of equal cardinality
  according to the sign of $\mu$, so a Möbius product $\prod_{d \dvd n} x_d^{\mu(n/d)}$ is a
  quotient of two products with the same number of factors; and
  $\sum_{d \dvd n} \mu(n/d) = 0$ for $n \ge 2$.
\end{lemma}
\begin{proof}
  \leanok
  Standard properties of $\mu$. The equal-cardinality statement is the binomial identity
  \[ \sum_{k} (-1)^k \binom{r}{k} = 0, \]
  where $r \ge 1$ is the number of prime divisors of $n'$.
\end{proof}

\section{Squares}

\begin{lemma}[Non-squares mod $m$]
  \label{lem:zmodsquare}
  \lean{ZMod.isSquare_neg_one_of_isSquare_div, ZMod.not_isSquare_neg_one_of_dvd,
    ZMod.not_isSquare_neg_one_of_four_dvd, ZMod.isUnit_intCast_of_isCoprime_of_dvd_add,
    ZMod.intCast_eq_neg_intCast_of_dvd_add}
  \leanok
  If $P \equiv -Q \bmod m$ with $Q$ a unit mod $m$, and $P/Q$ is a square in $\QQ$, then $-1$ is
  a square mod $m$. Moreover $-1$ is not a square mod $m$ whenever $m$ has a divisor
  $\equiv 3 \bmod 4$, and not mod $4$.
\end{lemma}
\begin{proof}
  \leanok
  Reduce a representation $P/Q = (u/v)^2$ modulo $m$ and use that $Q$ is invertible; the
  supplementary statements are the standard obstructions to $-1$ being a square.
\end{proof}

\begin{lemma}[Squares and coprimality]
  \label{lem:coprimesquare}
  \lean{Int.isSquare_abs_of_isSquare_prod_of_pairwise_isCoprime, prod_zpow_eq_zpow_sum}
  \leanok
  If a product of pairwise coprime integers is a square, then the absolute value of each factor
  is a square.
\end{lemma}
\begin{proof}
  \leanok
  Compare the $p$-adic valuations: each prime occurs in exactly one factor.
\end{proof}

\section{Even polynomials}

\begin{lemma}[Even polynomials are polynomials in $X^2$]
  \label{lem:evenpoly}
  \lean{Polynomial.even_eq_comp_X_sq_add_C, Polynomial.eq_expand_two_contract_of_comp_neg_X,
    Polynomial.expand_two_comp_neg_X, Polynomial.X_sq_add_C_comp_neg_X,
    Polynomial.comp_neg_X_eq_or_eq_neg_of_associated, Associated.comp_neg_X,
    Polynomial.normalize_normalize_comp_neg_X}
  \leanok
  Over a domain of characteristic $\ne 2$, a polynomial fixed by $X \mapsto -X$ is a polynomial
  in $X^2$, hence in $X^2 + c$ for any $c$. A polynomial merely \emph{associated} to its
  reflection satisfies $p(-X) = p$ or $p(-X) = -p$.
\end{lemma}
\begin{proof}
  \leanok
  The odd coefficients vanish because $2 \ne 0$; for the second statement, the unit relating $p$
  to $p(-X)$ is a constant squaring to $1$.
\end{proof}

\begin{lemma}[Pairing the irreducible factors]
  \label{lem:involution}
  \lean{Multiset.exists_add_map_of_involutive, normalizedFactors_map_mulEquiv_eq}
  \leanok
  A fixed-point-free involution of a multiset splits it as $N + \tau(N)$. Applied to
  $q \mapsto \mathrm{normalize}(q(-X))$ on the normalized factors of an even polynomial without
  nontrivial even divisors, this yields the factorization $\pm p = g \cdot g(-X)$.
\end{lemma}
\begin{proof}
  \leanok
  Induction on the multiset, removing an orbit $\{q, \tau q\}$ at a time; that the involution
  preserves the multiset of normalized factors is functoriality of factorization under the ring
  automorphism $X \mapsto -X$.
\end{proof}

\chapter{\texorpdfstring{Further values of $a$ with $\Om_n \cong \Ct{n}$: the results of Li}{Further
  values of a with the full wreath power: the results of Li}}
\label{chap:li}

This chapter formalizes the results of Li \cite{li2020, li2021} that extend
Theorem~\ref{thm:section3} to further values of $a$. Li's two papers use the notation of
Chapter~2 with the rescaled polynomial of Definition~\ref{def:normpoly}:
$g = |a| X^2 + \sgn(a)$, $\gamma_1 = 1$, $\gamma_{n+1} = g(\gamma_n)$, so that
$\gamma_n = |c_n|/|a|$ and $\beta_n = |b_n|$ for $n \ge 2$. (In \cite{li2020} the polynomial is
written $x^2 - a$ with $a > 0$; below, $a$ always denotes the coefficient of $f = X^2 + a$, so the
case of \cite{li2020} is $a < 0$, $a \equiv 3 \bmod 4$.) The Lean sources are in
\texttt{QuadraticIterates/Li/}.

\section{The results}

\begin{theorem}[\cite{li2021}, Theorem 3.3]
  \label{thm:li-A}
  \lean{QuadraticIterates.li2021_theorem_3_3, QuadraticIterates.nonempty_mulEquiv_of_forall_not_isSquare_abs_bSeq}
  \leanok
  \uses{def:galois, def:wreath}
  Let $a = -(8k+2)(8k+3)$ with $k \ge 0$. Then $\Om_n \cong \Ct{n}$ for all $n \ge 1$.
\end{theorem}
\begin{proof}
  \leanok
  \uses{lem:li-sqfree, lem:li-A-nonsqfree, thm:section1, def:normpoly}
  $-a = (8k+2)(8k+3) \equiv 6 \bmod 8$ is not a square. For $n \ge 2$, $|b_n| = \beta_n$ is not a
  square: by Lemma~\ref{lem:li-sqfree}~(2) if $n$ is squarefree, by
  Lemma~\ref{lem:li-A-nonsqfree} otherwise. Part~2) of Theorem~\ref{thm:section1}.
  \formalnote{The case split on \texttt{Squarefree n} in front of
  \lname{section1\_of\_not\_isSquare\_abs\_bSeq} is
  \lname{nonempty\_mulEquiv\_of\_forall\_not\_isSquare\_abs\_bSeq}, shared with
  Theorem~\ref{thm:li-B}.}
\end{proof}

\begin{theorem}[\cite{li2021}, Theorem 3.9]
  \label{thm:li-B}
  \lean{QuadraticIterates.li2021_theorem_3_9}
  \leanok
  \uses{def:galois, def:wreath}
  Let $a = -\bigl((4k+1)(4k+2) + 1\bigr)$ with $k \ge 0$. Then $\Om_n \cong \Ct{n}$ for all
  $n \ge 1$.
\end{theorem}
\begin{proof}
  \leanok
  \uses{lem:li-sqfree, lem:li-B-odd, lem:li-B-even, thm:section1, def:normpoly}
  $-a = (4k+1)(4k+2) + 1 \equiv 3 \bmod 4$ is not a square, and $a \equiv 1 \bmod 4$. For
  $n \ge 2$: if $n$ is squarefree, Lemma~\ref{lem:li-sqfree}~(3); if not and $4 \nmid n$,
  Lemma~\ref{lem:li-B-odd}; if $4 \dvd n$, then $n/\rad(n)$ is even and Lemma~\ref{lem:li-B-even}
  applies. Part~2) of Theorem~\ref{thm:section1}.
\end{proof}

\begin{theorem}[\cite{li2020}, Theorems 3.4 and 3.5]
  \label{thm:li-C}
  \lean{QuadraticIterates.li2020_theorem_3_5, QuadraticIterates.li2020_theorem_3_4, QuadraticIterates.bSeq_two}
  \leanok
  \uses{def:galois, def:wreath, def:bseq}
  Let $a < 0$ with $a \equiv 3 \bmod 4$ and $-a$ not a square. Then $|b_n|$ is not a square for
  all $n \ge 3$, and
  \[ \Om_n \cong \Ct{n} \text{ for all } n \ge 1 \iff -a - 1 \text{ is not a square in } \ZZ. \]
\end{theorem}
\begin{proof}
  \leanok
  \uses{lem:li-C-nonsqfree, lem:li-C-sqfree, thm:section1, def:normpoly, def:twoindep}
  For $n \ge 3$, $|b_n| = \beta_n$ is not a square: Lemma~\ref{lem:li-C-nonsqfree} if $n$ is not
  squarefree, Lemma~\ref{lem:li-C-sqfree} if it is (then $n > 2$). Moreover
  $b_2 = c_2/b_1 = -(a+1) = -a - 1 > 0$. If $-a - 1$ is not a square, no $|b_n|$ with $n \ge 2$
  is, and part~2) of Theorem~\ref{thm:section1} gives $\Om_n \cong \Ct{n}$ for all $n$. If
  $-a - 1 = b_2$ is a square, then $b_1, b_2$ are not $2$-independent, so by part~1) of
  Theorem~\ref{thm:section1} (the equivalence (a) $\Leftrightarrow$ (c) at $n = 2$),
  $\Om_2 \not\cong \Ct{2}$.
  \formalnote{The value $b_2 = -a - 1$ is \lname{bSeq\_two}, from $c_2 = b_1 b_2$
  (\lname{cSeq\_eq\_prod\_bSeq}). The direction ``$\Rightarrow$'' uses \lname{section1\_tfae}
  and \lname{TwoIndependent.not\_isSquare}. \cite{li2020} also shows that in the exceptional case
  $-a - 1 = m^2$ the index of $\Om_n$ in $\Ct{n}$ is exactly $2$ for all $n \ge 2$; this needs an
  analysis of the Kummer tower when the base is not maximal, which the present infrastructure
  (Lemma~\ref{lem:1.5} assumes $\Om_n \cong \Ct{n}$) does not provide, and it is not formalized.}
\end{proof}

The first two theorems cover $a = -6, -110, -342, \dots$ and $a = -3, -31, -91, \dots$; the third
covers $a = -13, -21, -29, \dots$ and identifies the failures $a = -5, -17, -37, \dots$ (where
$b_2 = -a - 1$ is a square). All three are proved, as Theorem~\ref{thm:section3} is, by showing
that no $|b_n|$ with $n \ge 2$ (resp.\ $n \ge 3$) is a square and invoking part~2) of
Theorem~\ref{thm:section1}; what changes is how the non-squareness of $\beta_n$ is established.
The three theorems are also part of the comparator setup of the repository, one challenge
statement each.

\emph{Numerical check.} All congruences and identities below were verified in PARI/GP for the
first members of each family and $n \le 13$ before the formalization, and the lists of $n$ with
$|b_n|$ a square were computed: empty for $a = -6, -110, -342, -3, -31, -91, -13, -21, -29, -33$
and $\{2\}$ for $a = -5, -17, -37$. The check exposed a misprint in \cite{li2021}, Lemma~2.2: its
third case must read $a \equiv 1 \bmod 4$, not $3 \bmod 4$ (the case $a \equiv 3 \bmod 4$ is the
one of \cite{li2020}, and $a = -5$ is a counterexample to the printed statement).

\section{The criterion of Lemma 2.1, one index at a time}

The proof of Lemma~\ref{lem:2.1} has three layers, which the formalization keeps apart because
Li's arguments reuse the two lower ones with different inputs: the arithmetic of the Möbius
product ($\beta_n$ as a quotient $P/Q$ of two products of the $\gamma_{kt}$, and the criterion
$P \equiv -Q \bmod m$), the congruence patterns that produce $P \equiv -Q$ (values of
$\gamma_{kt}$ depending only on a class of $t$, with a sign partition balanced on each class),
and the sequence facts feeding the patterns (eventual constancy on the multiples of $k$ modulo
$\gamma_k + \gamma_{2k}$; $2$-periodicity modulo $\gamma_k + \gamma_{k+2}$). The criteria of
\cite{li2021}, Section~2, are then four short lemmas (\ref{lem:li-crit}--\ref{lem:li-odd}), each
stated for a single index $n$ with its own modulus, and Lemmas~\ref{lem:2.1} and~\ref{lem:2.2}
are corollaries.

Throughout, $g$ is an even integer polynomial, $\varepsilon = \pm 1$, all $\gamma_n$ ($n \ge 1$)
are nonzero, and for $n \ge 2$ we write $n' := \rad(n)$ and $k := n/n'$; ``$m$ is a modulus''
means $m \ge 1$, and ``$x$ is a unit mod $m$'' means that $x$ and $m$ are coprime. Li always
takes $m$ to be a prime $p \equiv 3 \bmod 4$; the formalization keeps the modulus form of
Lemma~\ref{lem:2.1}, which needs no prime and turns ``$p \nmid \gamma$'' into coprimality. The
sign partition of the divisors $t$ of $n'$ is the pair of sets
$S_\pm := \{t \dvd n' : \mu(n'/t) = \pm 1\}$ (Lemma~\ref{lem:moebiussign}).

\begin{lemma}[Balanced classes]
  \label{lem:li-classes}
  \lean{Finset.prod_eq_neg_prod_of_forall_card_filter_eq, Finset.card_filter_eq_one_eq_card_filter_eq_neg_one_of_sum_eq_zero}
  \leanok
  Let $R$ be a commutative ring, $S_+$ and $S_-$ disjoint finite sets with $a \in S_+ \cup S_-$,
  $x \colon S_+ \cup S_- \to R$, $w \colon I \to R$ and $\kappa \colon S_+ \cup S_- \to I$ a
  ``class'' function such that $x_t = w_{\kappa(t)}$ for $t \ne a$ and $x_a = -w_{\kappa(a)}$. If
  every class is balanced, i.e.\ $\#\{t \in S_+ : \kappa(t) = i\} = \#\{t \in S_- : \kappa(t) = i\}$
  for all $i$, then
  \[ \prod_{t \in S_+} x_t = -\prod_{t \in S_-} x_t. \]
\end{lemma}
\begin{proof}
  \leanok
  Group each product by classes (\lname{Finset.prod\_comp}): on the class $i$ the product over
  $S_\pm$ is $w_i^{\#\{t \in S_\pm : \kappa(t) = i\}}$, times $-1$ if $a \in S_\pm$. The exponents
  agree by balance, and exactly one of $S_+$, $S_-$ contains $a$.
  \formalnote{With $I$ a one-point set this gives \lname{prod\_gammaSeq\_mul\_eq\_neg\_prod}
  (Lemma~\ref{lem:li-crit}), with $I = \ZZ/2$ the parity classes of Lemma~\ref{lem:li-odd}. The
  companion \lname{Finset.card\_filter\_eq\_one\_eq\_card\_filter\_eq\_neg\_one\_of\_sum\_eq\_zero}
  produces the balance: a $\pm 1$-valued function summing to $0$ on a finite set takes each sign
  equally often.}
\end{proof}

\begin{lemma}[Parity-balanced sign partition]
  \label{lem:li-parity}
  \lean{ArithmeticFunction.sum_divisors_filter_mod_two_moebius_div, ArithmeticFunction.sum_divisors_filter_dvd_moebius_div, Squarefree.card_filter_moebius_div_eq_one_eq_card_filter_eq_neg_one}
  \leanok
  \uses{lem:moebiussign}
  Let $n \ge 3$. Then
  \[ \sum_{t \dvd n,\ t \text{ odd}} \mu(n/t) = 0
     = \sum_{t \dvd n,\ t \text{ even}} \mu(n/t). \]
  Consequently, for squarefree $n' \ge 3$, the sign partition $S_+ \sqcup S_-$ of the divisors of
  $n'$ is balanced on the odd and on the even divisors separately.
\end{lemma}
\begin{proof}
  \leanok
  The sum over the even divisors is the restricted Möbius sum $\sum_{t \dvd n,\ 2 \dvd t} \mu(n/t)$,
  which is $[n = 2] = 0$ (Lemma~\ref{lem:moebiussign},
  \lname{sum\_divisors\_filter\_dvd\_moebius\_div}); the sum over the odd divisors is the
  difference of the total sum $\sum_{t \dvd n} \mu(n/t) = 0$ and the even one. For squarefree
  $n'$ the values $\mu(n'/t)$ are $\pm 1$, so on any set of divisors with vanishing sum the two
  signs occur equally often (\lname{Squarefree.card\_filter\_moebius\_div\_eq\_one\_eq\_card\_filter\_eq\_neg\_one}).
  \formalnote{The hypothesis $n \ge 3$ only excludes $n = 2$, where the odd class $\{1\}$ is
  unbalanced; no squarefreeness is needed for the vanishing of the sums. The residue classes
  are written as $t \bmod 2 = i$ with $i$ ranging over $\NN$ (the classes $i \ge 2$ are empty),
  which is the form the class function $\kappa(t) = t \bmod 2$ of Lemma~\ref{lem:li-classes}
  consumes.}
\end{proof}

\begin{lemma}[$\beta_n$ from a numerator/denominator congruence]
  \label{lem:li-PQ}
  \lean{QuadraticIterates.not_isSquare_betaSeq_of_prod_eq_neg_prod}
  \leanok
  \uses{def:gammaseq, lem:moebiussign, lem:zmodsquare}
  Let $n \ge 1$, $n' := \rad(n)$, $k := n/n'$, and put $P := \prod_{t \in S_+} \gamma_{kt}$,
  $Q := \prod_{t \in S_-} \gamma_{kt}$. Suppose $P = c P'$ and $Q = c Q'$ with an integer
  $c \ne 0$, and let $m$ be a modulus with $-1$ not a square mod $m$ such that
  $P' \equiv -Q' \bmod m$ and $Q'$ is a unit mod $m$. Then $\beta_n$ is not a square in $\QQ$.
\end{lemma}
\begin{proof}
  \leanok
  $\beta_n = \prod_{d \dvd n} \gamma_d^{\mu(n/d)} = P/Q = P'/Q'$ by
  Lemma~\ref{lem:moebiussign} (\lname{prod\_pow\_moebius\_eq\_div}), and a quotient of two
  integers with $P' \equiv -Q'$, $Q'$ a unit, is not a rational square when $-1$ is not a square
  mod $m$ (Lemma~\ref{lem:zmodsquare}, \lname{ZMod.isSquare\_neg\_one\_of\_isSquare\_div}).
  \formalnote{The common factor $c$ is $1$ everywhere except in Lemma~\ref{lem:li-C-sqfree},
  where it is a power of $\gamma_2$ (Li works modulo a higher power of the modulus instead). The
  hypothesis ``$P' \equiv -Q'$ with $Q'$ a unit'' is stated in $\ZZ/m\ZZ$.}
\end{proof}

\begin{lemma}[\cite{li2021}, Lemma 2.1: Stoll's criterion at one index]
  \label{lem:li-crit}
  \lean{QuadraticIterates.not_isSquare_betaSeq_of_dvd_add_two_mul, QuadraticIterates.prod_gammaSeq_mul_eq_neg_prod, QuadraticIterates.isUnit_prod_gammaSeq_mul}
  \leanok
  \uses{def:gammaseq}
  Let $n \ge 2$, $k := n/\rad(n)$, and let $m$ be a modulus with $m \dvd \gamma_k + \gamma_{2k}$,
  $m$ coprime to $\gamma_k$, and $-1$ not a square mod $m$. Then $\beta_n$ is not a square in
  $\QQ$.
\end{lemma}
\begin{proof}
  \leanok
  \uses{lem:li-classes, lem:li-PQ, lem:moebiussign}
  Exactly the proof of Lemma~\ref{lem:2.1}: mod $m$, $\gamma_{tk} \equiv \gamma_{2k}$ for all
  $t \ge 2$ (\lname{gammaSeq\_mul\_eq\_two\_mul}) and $\gamma_k \equiv -\gamma_{2k}$, so
  Lemma~\ref{lem:li-classes} with a single class gives $P \equiv -Q$
  (\lname{prod\_gammaSeq\_mul\_eq\_neg\_prod}), $Q$ is a unit because
  $\gamma_{2k} \equiv -\gamma_k$ is (\lname{isUnit\_prod\_gammaSeq\_mul}), and
  Lemma~\ref{lem:li-PQ} concludes.
  \formalnote{Lemma~\ref{lem:2.1} (\lname{not\_isSquare\_betaSeq}) is recovered by instantiating
  the modulus at the index $k = n/\rad(n)$.}
\end{proof}

\begin{lemma}[\cite{li2021}, Corollary 2.4]
  \label{lem:li-succ}
  \lean{QuadraticIterates.not_isSquare_betaSeq_of_dvd_add_succ, QuadraticIterates.not_isSquare_abs_bSeq_of_dvd_add_succ}
  \leanok
  \uses{def:gammaseq}
  Let $g(0)$ be a unit, $n \ge 2$, $k := n/\rad(n)$, and $m$ a modulus with
  $m \dvd \gamma_k + \gamma_{k+1}$ and $-1$ not a square mod $m$. Then $\beta_n$ is not a square
  in $\QQ$.
\end{lemma}
\begin{proof}
  \leanok
  \uses{lem:li-crit}
  $\gamma_k + \gamma_{k+1}$ divides $\gamma_k + \gamma_{2k}$ (\lname{gammaSeq\_add\_succ\_dvd})
  and is coprime to $\gamma_k$ (\lname{isCoprime\_gammaSeq\_add\_succ}); both properties pass to
  the divisor $m$, and Lemma~\ref{lem:li-crit} applies.
  \formalnote{Lemma~\ref{lem:2.2} goes through this statement
  (\lname{not\_isSquare\_betaSeq\_of\_pos\_of\_not\_isSquare\_neg\_one} is the case
  $m = \gamma_k + \gamma_{k+1}$). The per-index form is what Theorem~\ref{thm:li-C} needs: there
  $\gamma_1 + \gamma_2 = -a \equiv 1 \bmod 4$ fails the hypothesis, but every
  $\gamma_k + \gamma_{k+1}$ with $k \ge 2$ satisfies it.}
\end{proof}

\begin{lemma}[$2$-periodicity]
  \label{lem:li-period}
  \lean{QuadraticIterates.gammaSeq_eq_ite_even_of_add_two_eq_zero, QuadraticIterates.gammaSeq_mul_eq_ite_even_of_odd, QuadraticIterates.gammaSeq_eq_ite_even_of_add_two_eq}
  \leanok
  \uses{def:gammaseq}
  Over any commutative ring, if $k \ge 1$ and $\gamma_k + \gamma_{k+2} = 0$, then for all
  $r \ge k + 1$
  \[ \gamma_r = \begin{cases} -\gamma_k & \text{if } r \equiv k \bmod 2, \\
                              \gamma_{k+1} & \text{otherwise.} \end{cases} \]
  In particular, for odd $k$ and $t \ge 2$, $\gamma_{kt} = -\gamma_k$ if $t$ is odd and
  $\gamma_{kt} = \gamma_{k+1}$ if $t$ is even.
\end{lemma}
\begin{proof}
  \leanok
  $g(-\gamma_k) = g(\gamma_k) = \gamma_{k+1}$ since $g$ is even, and
  $g(\gamma_{k+1}) = \gamma_{k+2} = -\gamma_k$; so the sequence alternates between the two values
  from index $k + 1$ on. For odd $k$, $kt \equiv k \bmod 2$ exactly when $t$ is odd.
  \formalnote{The same induction, without evenness, gives the $2$-periodicity from a $2$-cycle
  $\gamma_{n_0+2} = \gamma_{n_0}$ (\lname{gammaSeq\_eq\_ite\_even\_of\_add\_two\_eq}), used for
  the residue computations of Section~\ref{sec:li-residues}. All are applied in $\ZZ/m\ZZ$
  through \lname{intCast\_gammaSeq}, as the existing congruence lemmas are.}
\end{proof}

\begin{proposition}[\cite{li2021}, Proposition 2.5, $k$ even]
  \label{lem:li-even}
  \lean{QuadraticIterates.not_isSquare_betaSeq_of_even_of_dvd_add_two, QuadraticIterates.not_isSquare_abs_bSeq_of_even_of_dvd_add_two, isCoprime_of_isCoprime_gcd_of_dvd_add, IsRelPrime.of_gcd_right_of_dvd_add}
  \leanok
  \uses{def:gammaseq}
  Let $n \ge 2$ with $k = n/\rad(n)$ even, and $m$ a modulus with $m \dvd \gamma_k + \gamma_{k+2}$,
  $m$ coprime to $\gamma_2$, and $-1$ not a square mod $m$. Then $\beta_n$ is not a square.
\end{proposition}
\begin{proof}
  \leanok
  \uses{lem:li-period, lem:li-crit, lem:strongdvd}
  Mod $m$, Lemma~\ref{lem:li-period} gives $\gamma_{2k} \equiv -\gamma_k$ ($2k \equiv k \bmod 2$),
  so $m \dvd \gamma_k + \gamma_{2k}$. Strong divisibility (\lname{gammaSeq\_associated\_gcd},
  $\gcd(k, k+2) = 2$) gives $\gcd(\gamma_k, \gamma_{k+2}) = \gamma_2$; a common divisor of $m$
  and $\gamma_k$ divides $\gamma_{k+2}$, hence $\gamma_2$, hence is $1$
  (\lname{isCoprime\_of\_isCoprime\_gcd\_of\_dvd\_add}, over any Bézout domain).
  Lemma~\ref{lem:li-crit}.
\end{proof}

\begin{proposition}[\cite{li2021}, Proposition 2.5, $k$ odd]
  \label{lem:li-odd}
  \lean{QuadraticIterates.not_isSquare_betaSeq_of_odd_of_dvd_add_two, QuadraticIterates.not_isSquare_abs_bSeq_of_odd_of_dvd_add_two, QuadraticIterates.prod_gammaSeq_mul_eq_neg_prod_of_odd, QuadraticIterates.isUnit_prod_gammaSeq_mul_of_odd, QuadraticIterates.isCoprime_gammaSeq_of_odd_of_dvd_add_two}
  \leanok
  \uses{def:gammaseq}
  Let $g(0)$ be a unit, $n \ge 2$ with $k = n/\rad(n)$ odd and $k > 1$, and $m$ a modulus with
  $m \dvd \gamma_k + \gamma_{k+2}$, $m$ coprime to $\gamma_{k+1}$, and $-1$ not a square mod $m$.
  Then $\beta_n$ is not a square.
\end{proposition}
\begin{proof}
  \leanok
  \uses{lem:li-period, lem:li-parity, lem:li-classes, lem:li-PQ, lem:strongdvd}
  Since $k$ is odd and $> 1$, $n' = \rad(n)$ is neither $1$ nor $2$ (an odd prime factor of $k$
  divides $n'$). Mod $m$, by Lemma~\ref{lem:li-period}, $\gamma_{kt} \equiv -\gamma_k$ for odd
  $t \ge 3$ and $\gamma_{kt} \equiv \gamma_{k+1}$ for even $t$. Both values are units mod $m$:
  $\gamma_{k+1}$ by hypothesis, $\gamma_k$ because $\gcd(\gamma_k, \gamma_{k+2}) = \gamma_1$ is a
  unit ($\gcd(k, k+2) = 1$), so a common divisor of $m$ and $\gamma_k$ divides $\gamma_{k+2}$ and
  hence $\gamma_1$ (\lname{isCoprime\_gammaSeq\_of\_odd\_of\_dvd\_add\_two}). Take the
  parity-balanced partition of Lemma~\ref{lem:li-parity} and apply Lemma~\ref{lem:li-classes}
  with the two classes ``odd'' ($w = -\gamma_k$; at $t = 1$ the value $\gamma_k = -w$) and
  ``even'' ($w = \gamma_{k+1}$): $P \equiv -Q$ with $Q$ a unit
  (\lname{prod\_gammaSeq\_mul\_eq\_neg\_prod\_of\_odd},
  \lname{isUnit\_prod\_gammaSeq\_mul\_of\_odd}). Lemma~\ref{lem:li-PQ}.
  \formalnote{Li derives the coprimality of $p$ and $\gamma_{k+1}$ from $p$ being an odd prime;
  in modulus form it is a hypothesis, discharged for the rescaled polynomial by
  Lemma~\ref{lem:li-cop}.}
\end{proof}

\section{Residues of the rescaled sequence}
\label{sec:li-residues}

In this section $a < 0$, $g = |a| X^2 - 1$ is the rescaled polynomial of
Definition~\ref{def:normpoly}, so $\gamma_1 = 1$, $\gamma_2 = |a| - 1$, and all $\gamma_n$ are
positive. The residue computations are done in $\ZZ/8\ZZ$ or $\ZZ/4\ZZ$ by the $2$-periodicity
of Lemma~\ref{lem:li-period}.

\begin{lemma}[The sequence modulo $\gamma_2^2$]
  \label{lem:li-modsq}
  \lean{QuadraticIterates.gammaSeq_two_sq_dvd_sub_ite_even, QuadraticIterates.gammaSeq_two_dvd_of_even}
  \leanok
  \uses{def:normpoly}
  For every even $g$ with $g(0)^2 = 1$ and $\gamma_1 = 1$, and all $n \ge 2$:
  $\gamma_n \equiv \gamma_2$ for even $n$ and $\gamma_n \equiv g(0) \bmod \gamma_2^2$ for odd $n$.
\end{lemma}
\begin{proof}
  \leanok
  By induction. If $\gamma_n \equiv \gamma_2$ then $\gamma_2 \dvd \gamma_n$, so
  $\gamma_{n+1} = g(\gamma_n) \equiv g(0)$ because $\gamma_n^2 \dvd g(\gamma_n) - g(0)$
  (\lname{sq\_dvd\_gammaSeq\_succ\_sub}); if $\gamma_n \equiv g(0)$ then
  $\gamma_n^2 \equiv g(0)^2 = 1$, so $\gamma_{n+1} = g(\gamma_n) \equiv g(1) = \gamma_2$
  (\lname{EvenPoly.dvd\_eval\_sub}).
  \formalnote{Used mod $\gamma_2$ in Lemma~\ref{lem:li-C-sqfree} and mod $(4k+1)^2$ in
  Lemma~\ref{lem:li-B-even}.}
\end{proof}

\begin{lemma}[$a$ even]
  \label{lem:li-mod8-2}
  \lean{QuadraticIterates.gammaSeq_normPoly_zmod_eight_of_even}
  \leanok
  \uses{def:normpoly}
  If $a < 0$ is even, then $\gamma_n \equiv -a - 1 \bmod 8$ for all $n \ge 2$; in particular
  $\gamma_n \equiv 5 \bmod 8$ when $a \equiv 2 \bmod 8$.
\end{lemma}
\begin{proof}
  \leanok
  In $\ZZ/8\ZZ$, $\gamma_2 = g(1) = -a - 1$ is odd, so $\gamma_2^2 = 1$ and $g(\gamma_2) = g(1)$:
  a fixed point (\lname{gammaSeq\_eq\_eval\_one}).
\end{proof}

\begin{lemma}[$a \equiv 1 \bmod 4$]
  \label{lem:li-mod8-1}
  \lean{QuadraticIterates.gammaSeq_normPoly_zmod_eight_of_emod_four_eq_one, QuadraticIterates.gammaSeq_normPoly_add_two_emod_eight_eq_six, QuadraticIterates.gammaSeq_normPoly_emod_four_eq_two_of_even}
  \leanok
  \uses{def:normpoly}
  If $a < 0$ and $a \equiv 1 \bmod 4$, then $\gamma_n \equiv 3 \bmod 8$ for odd $n \ge 3$ and
  $\gamma_n \equiv -a - 1 \bmod 8$ (which is $2$ or $6$) for even $n \ge 2$. Consequently
  $\gamma_k + \gamma_{k+2} \equiv 6 \bmod 8$ for odd $k \ge 3$, and $\gamma_k \equiv 2 \bmod 4$ for
  even $k \ge 2$.
\end{lemma}
\begin{proof}
  \leanok
  \uses{lem:li-period}
  In $\ZZ/8\ZZ$, $|a| \in \{3, 7\}$, $\gamma_2 = |a| - 1 \in \{2, 6\}$, $\gamma_3 = |a| \cdot 4 - 1
  = 3$, $\gamma_4 = |a| \cdot 9 - 1 = |a| - 1 = \gamma_2$, and the pattern repeats with period $2$
  (Lemma~\ref{lem:li-period}, $2$-cycle form). The two evaluations are \texttt{decide} after a
  case split on $|a| \bmod 8$.
\end{proof}

\begin{lemma}[$a \equiv 3 \bmod 4$]
  \label{lem:li-mod4-3}
  \lean{QuadraticIterates.gammaSeq_normPoly_zmod_four_of_emod_four_eq_three, QuadraticIterates.gammaSeq_normPoly_add_succ_emod_four_eq_three, QuadraticIterates.gammaSeq_eq_ite_even_of_eval_one_eq_zero}
  \leanok
  \uses{def:normpoly}
  If $a < 0$ and $a \equiv 3 \bmod 4$, then $\gamma_n \equiv 0 \bmod 4$ for even $n \ge 2$ and
  $\gamma_n \equiv 3 \bmod 4$ for odd $n \ge 3$; hence $\gamma_k + \gamma_{k+1} \equiv 3 \bmod 4$
  for all $k \ge 2$.
\end{lemma}
\begin{proof}
  \leanok
  \uses{lem:li-period}
  In $\ZZ/4\ZZ$, $g = X^2 - 1$, $\gamma_2 = 0$, $\gamma_3 = -1$, $\gamma_4 = 0$, and the pattern
  repeats.
  \formalnote{The pattern $1, 0, g(0), 0, g(0), \dots$ holds over any ring for an even $g$ with
  $g(1) = 0$, $g(0)^2 = 1$ and $\gamma_1 = 1$ (\lname{gammaSeq\_eq\_ite\_even\_of\_eval\_one\_eq\_zero});
  the same lemma is the core of Lemma~\ref{lem:li-cop}.}
\end{proof}

\begin{lemma}[Coprimality at odd indices]
  \label{lem:li-cop}
  \lean{QuadraticIterates.dvd_two_of_dvd_add_two_of_dvd_succ, QuadraticIterates.isCoprime_gammaSeq_succ_of_odd_of_add_two_eq_two_mul}
  \leanok
  \uses{def:normpoly}
  Let $g$ be even with $g(0)^2 = 1$ and $\gamma_1 = 1$, and let $k \ge 3$ be odd. Then every
  common divisor of $\gamma_k + \gamma_{k+2}$ and $\gamma_{k+1}$ divides $2$; in particular, if
  $\gamma_k + \gamma_{k+2} = 2M$ with $M$ odd, then $M$ is coprime to $\gamma_{k+1}$.
\end{lemma}
\begin{proof}
  \leanok
  \uses{lem:li-mod4-3}
  Let $d$ be a common divisor and compute in $\ZZ/d\ZZ$: $\gamma_{k+1} = 0$, so
  $\gamma_{k+2} = g(0)$ and $\gamma_k = -g(0)$; then
  $0 = \gamma_{k+1} = g(\gamma_k) = g(-g(0)) = g(g(0)) = g(1)$ (evenness and $g(0)^2 = 1$). For
  such a $g$ the sequence is $1, 0, g(0), 0, g(0), \dots$ (Lemma~\ref{lem:li-mod4-3}, general
  form), so $\gamma_k = g(0)$ for the odd $k \ge 3$; together with $\gamma_k = -g(0)$ this gives
  $2 = 0$ in $\ZZ/d\ZZ$.
  \formalnote{Applied to $g = |a| X^2 - 1$ with $a \equiv 1 \bmod 4$ in Lemma~\ref{lem:li-B-odd},
  where $\gamma_k + \gamma_{k+2} \equiv 6 \bmod 8$ makes $M \equiv 3 \bmod 4$.}
\end{proof}

\section{The applications}

\begin{lemma}[\cite{li2021}, Lemma 2.2: squarefree indices]
  \label{lem:li-sqfree}
  \lean{QuadraticIterates.not_isSquare_abs_bSeq_of_squarefree, QuadraticIterates.not_isSquare_abs_bSeq_of_squarefree_of_pos_of_emod_eight_eq_four, QuadraticIterates.not_isSquare_abs_bSeq_of_squarefree_of_neg_of_emod_eight_eq_two, QuadraticIterates.not_isSquare_abs_bSeq_of_squarefree_of_neg_of_emod_four_eq_one}
  \leanok
  \uses{def:normpoly}
  Let $n > 1$ be squarefree. Then $|b_n|$ is not a square as soon as $-1$ is not a square modulo
  $\gamma_1 + \gamma_2 = |a| + \sgn(a) + 1$; this holds in each of the cases
  (1) $a > 0$, $a \equiv 4 \bmod 8$; (2) $a < 0$, $a \equiv 2 \bmod 8$; (3) $a < 0$,
  $a \equiv 1 \bmod 4$.
\end{lemma}
\begin{proof}
  \leanok
  \uses{lem:li-succ}
  Here $k = 1$, so Lemma~\ref{lem:li-succ} with $m := \gamma_1 + \gamma_2 = 1 + g(1)$ applies as
  soon as $-1$ is not a square mod $m$: $m = a + 2 \equiv 6 \bmod 8$ in case (1), $m = -a \equiv
  6 \bmod 8$ in case (2), and $m = -a \equiv 3 \bmod 4$ in case (3)
  (\lname{ZMod.not\_isSquare\_neg\_one\_of\_emod\_eight\_eq\_six},
  \lname{\dots\_of\_emod\_four\_eq\_three}).
  \formalnote{Case (3) is the misprinted one, see above. Stoll's cases $a > 0$,
  $a \equiv 1, 2 \bmod 4$ and $a < 0$, $a \equiv 0 \bmod 4$ fit the same statement; they are not
  repeated here because Lemma~\ref{lem:2.2} already gives them for all $n$.}
\end{proof}

\subsection*{Theorem A}

\begin{lemma}[\cite{li2021}, Lemma 3.1]
  \label{lem:li-A-nonsqfree}
  \lean{QuadraticIterates.not_isSquare_abs_bSeq_of_not_squarefree_of_eq_neg_mul}
  \leanok
  \uses{def:normpoly}
  Let $a = -(8k+2)(8k+3)$ and let $n \ge 1$ be not squarefree. Then $|b_n|$ is not a square.
\end{lemma}
\begin{proof}
  \leanok
  \uses{lem:li-succ, lem:li-mod8-2}
  Let $k_0 := n/\rad(n) \ge 2$. Since $|a|\,y^2 + y - 1 = \bigl((8k+3)y - 1\bigr)\bigl((8k+2)y +
  1\bigr)$, the number $m := (8k+2)\gamma_{k_0} + 1$ divides
  $\gamma_{k_0} + \gamma_{k_0+1} = |a|\gamma_{k_0}^2 + \gamma_{k_0} - 1$, and
  $m \equiv 2 \cdot 5 + 1 \equiv 3 \bmod 4$ by Lemma~\ref{lem:li-mod8-2}. Lemma~\ref{lem:li-succ}.
  \formalnote{Here $a \equiv 2 \bmod 8$ because $(8k+2)(8k+3) \equiv 6 \bmod 8$. The
  discriminant $1 - 4a = (16k+5)^2$ is what makes the quadratic in $y$ split; \cite{li2021},
  Remark~3.2, notes that the complementary family $-(8k+5)(8k+6)$ is out of reach of this
  method ($a = -30$: no prime $\equiv 3 \bmod 4$ divides $\gamma_2 + \gamma_3$).}
\end{proof}

\subsection*{Theorem B}

\begin{lemma}[\cite{li2021}, Lemma 3.4]
  \label{lem:li-B-odd}
  \lean{QuadraticIterates.not_isSquare_abs_bSeq_of_not_squarefree_of_not_four_dvd_of_emod_four_eq_one, Nat.even_div_radical_iff_four_dvd, Nat.two_le_div_radical_of_not_squarefree, Nat.squarefree_iff_radical_eq_self}
  \leanok
  \uses{def:normpoly}
  Let $a < 0$, $a \equiv 1 \bmod 4$, and let $n \ge 1$ be not squarefree with $4 \nmid n$. Then
  $|b_n|$ is not a square.
\end{lemma}
\begin{proof}
  \leanok
  \uses{lem:li-odd, lem:li-mod8-1, lem:li-cop}
  $k = n/\rad(n)$ is odd (as $4 \nmid n$; \lname{Nat.even\_div\_radical\_iff\_four\_dvd}) and
  $> 1$ (as $n$ is not squarefree), so $k \ge 3$. By Lemma~\ref{lem:li-mod8-1},
  $\gamma_k + \gamma_{k+2} \equiv 6 \bmod 8$, so $M := (\gamma_k + \gamma_{k+2})/2 \equiv 3 \bmod 4$;
  it is coprime to $\gamma_{k+1}$ by Lemma~\ref{lem:li-cop}. Proposition~\ref{lem:li-odd}.
\end{proof}

\begin{lemma}[\cite{li2021}, Lemma 3.6]
  \label{lem:li-B-even}
  \lean{QuadraticIterates.not_isSquare_abs_bSeq_of_even_div_radical_of_eq_neg_mul_add_one}
  \leanok
  \uses{def:normpoly}
  Let $a = -\bigl((4k+1)(4k+2) + 1\bigr)$, $A := 4k+1$, $B := 4k+2$, $\alpha := -a = AB + 1$, and
  let $n \ge 1$ be such that $k_0 := n/\rad(n)$ is even. Then $|b_n|$ is not a square.
\end{lemma}
\begin{proof}
  \leanok
  \uses{lem:li-even, lem:li-mod8-1, lem:li-modsq, lem:strongdvd}
  \emph{The factorization.} In $\ZZ[y]$, with $g(y) = \alpha y^2 - 1$,
  \[ y + g(g(y))
     = (\alpha y + A)\,\bigl(\alpha^2 y^3 - A\alpha\, y^2 - (\alpha + B)\, y + B\bigr) \]
  (a polynomial identity in $k$ and $y$; the quotient was found by polynomial division in
  PARI). Hence $\alpha\gamma_{k_0} + A \dvd \gamma_{k_0} + \gamma_{k_0+2}$.

  \emph{The modulus.} $\gamma_2 = AB$ divides $\gamma_{k_0}$ (strong divisibility, $2 \dvd k_0$),
  and $\alpha \equiv 1 \bmod A$, so $A \dvd \alpha\gamma_{k_0} + A$; put
  $m := (\alpha\gamma_{k_0} + A)/A$. Then $m \equiv 3 \bmod 4$: $\gamma_{k_0} \equiv 2 \bmod 4$
  (Lemma~\ref{lem:li-mod8-1}), $\alpha \equiv 3 \bmod 4$, $A \equiv 1 \bmod 4$.

  \emph{Coprimality with $\gamma_2 = AB$.} Modulo $A^2$: $\gamma_{k_0} \equiv \gamma_2 = AB
  \equiv A \bmod A^2$ (Lemma~\ref{lem:li-modsq} reduced from $\gamma_2^2$ to $A^2$, and
  $AB = A^2 + A$), so $\alpha\gamma_{k_0} + A \equiv (\alpha + 1)A$ and $m \equiv \alpha + 1 \equiv
  2 \bmod A$, coprime to the odd $A$. Modulo $B = 2(2k+1)$: $m$ is odd, and since
  $B \dvd \gamma_{k_0}$, $A(m - 1) = \alpha\gamma_{k_0} \equiv 0 \bmod (2k+1)$ with
  $A \equiv -1 \bmod (2k+1)$, so $m \equiv 1 \bmod (2k+1)$.

  Proposition~\ref{lem:li-even} with this $m$.
  \formalnote{The discriminant $4\alpha - 3 = (8k+3)^2$ is what makes $y + g(g(y))$ split off a
  linear factor; \cite{li2021}, Remark~3.8, records that $-((4k+2)(4k+3)+1)$ (e.g.\ $a = -507$)
  is out of reach. The lemma covers $k_0 = 2$ as well, so \cite{li2021}, Lemma~3.5 (the case
  $k_0 = 2$ for arbitrary $a \equiv 1 \bmod 4$, with the modulus $(\gamma_2 + \gamma_4)/(2\gamma_2)$)
  is not needed and is not formalized; note that its printed claim
  $\gcd((\gamma_2+\gamma_4)/\gamma_2, \gamma_2) = 1$ is false (the gcd is $2$, e.g.\ $a = -3$:
  $182$ and $2$), and the odd part is what the argument uses. The Lean proof follows the three
  paragraphs above: the factorization is \lname{add\_eval\_eval\_eq\_mul}, the modulus with
  its residue is \lname{exists\_mul\_eq\_and\_emod\_four\_eq\_three\_of\_even}, its
  divisibility property \lname{dvd\_add\_two\_of\_mul\_eq}, and the coprimality argument is
  the Bézout-style \lname{isCoprime\_of\_mul\_eq}, applied in
  \lname{isCoprime\_gammaSeq\_two\_of\_mul\_eq}.}
\end{proof}

\subsection*{Theorem C}

\begin{lemma}[\cite{li2020}, Lemma 3.1]
  \label{lem:li-C-nonsqfree}
  \lean{QuadraticIterates.not_isSquare_abs_bSeq_of_not_squarefree_of_neg_of_emod_four_eq_three}
  \leanok
  \uses{def:normpoly}
  Let $a < 0$, $a \equiv 3 \bmod 4$, and let $n \ge 1$ be not squarefree. Then $|b_n|$ is not a
  square.
\end{lemma}
\begin{proof}
  \leanok
  \uses{lem:li-succ, lem:li-mod4-3}
  $k = n/\rad(n) \ge 2$, so $\gamma_k + \gamma_{k+1} \equiv 3 \bmod 4$ (Lemma~\ref{lem:li-mod4-3});
  Lemma~\ref{lem:li-succ} with $m := \gamma_k + \gamma_{k+1}$.
\end{proof}

\begin{lemma}[\cite{li2020}, Lemmas 3.2 and 3.3]
  \label{lem:li-C-sqfree}
  \lean{QuadraticIterates.not_isSquare_abs_bSeq_of_squarefree_of_neg_of_emod_four_eq_three, QuadraticIterates.not_isSquare_betaSeq_of_squarefree_of_not_isSquare_neg_one}
  \leanok
  \uses{def:normpoly}
  Let $a < 0$, $a \equiv 3 \bmod 4$, and let $n > 2$ be squarefree. Then $|b_n|$ is not a square.
\end{lemma}
\begin{proof}
  \leanok
  \uses{lem:li-classes, lem:li-parity, lem:li-PQ, lem:li-modsq}
  Here $k = 1$; take $m := \gamma_2 = |a| - 1$, which is $\equiv 0 \bmod 4$, so $-1$ is not a
  square mod $m$ (\lname{ZMod.not\_isSquare\_neg\_one\_of\_four\_dvd}). By
  Lemma~\ref{lem:li-modsq}, every even divisor $t$ of $n$ has $\gamma_t = m\,u_t$ with
  $u_t \equiv 1 \bmod m$, and every odd divisor $t \ge 3$ has $\gamma_t \equiv -1 \bmod m$;
  $\gamma_1 = 1$. With the parity-balanced partition of Lemma~\ref{lem:li-parity} ($n > 2$),
  $P = m^N P'$ and $Q = m^N Q'$ with $N$ the common number of even divisors in $S_+$ and $S_-$,
  and $P'$, $Q'$ the products of the $\gamma_t$ ($t$ odd) and $u_t$ ($t$ even).
  Lemma~\ref{lem:li-classes} with the classes ``odd'' ($w = -1$) and ``even'' ($w = 1$) gives
  $P' \equiv -Q'$ with $Q' \equiv \pm 1$; Lemma~\ref{lem:li-PQ} with $c = m^N$.
  \formalnote{Li treats odd $n$ (all divisors odd; his Lemma~3.2, modulo $4$) and even $n$
  (Lemma~3.3, modulo a power of $\gamma_2$) separately; modulo $\gamma_2$ the argument is the
  same for both, which is why the two lemmas are one declaration. The general statement
  (\lname{not\_isSquare\_betaSeq\_of\_squarefree\_of\_not\_isSquare\_neg\_one}) holds for every
  even $g$ with $g(0) = -1$ and $\gamma_1 = 1$, with the modulus $\gamma_2$.}
\end{proof}


\begin{thebibliography}{9}
  \bibitem{cremona}
    J.~E.~Cremona, \emph{On the Galois groups of the iterates of $x^2 + 1$}.
    Mathematika \textbf{36}, 259--261 (1989).
  \bibitem{li2020}
    H.-C.~Li, \emph{Arboreal Galois representation for a certain type of quadratic polynomials}.
    Arch.\ Math.\ \textbf{114}, 265--269 (2020).
    \href{https://doi.org/10.1007/s00013-019-01390-x}{DOI: 10.1007/s00013-019-01390-x}.
  \bibitem{li2021}
    H.-C.~Li, \emph{On Stoll's criterion for the maximality of quadratic arboreal Galois
    representations}. Arch.\ Math.\ \textbf{117}, 133--140 (2021).
    \href{https://doi.org/10.1007/s00013-021-01609-w}{DOI: 10.1007/s00013-021-01609-w}.
  \bibitem{odoni}
    R.~W.~K.~Odoni, \emph{Realising wreath products of cyclic groups as Galois groups}.
    Mathematika \textbf{35}, 101--113 (1988).
  \bibitem{paper}
    M.~Stoll, \emph{Galois groups over $\QQ$ of some iterated polynomials}.
    Arch.\ Math.\ \textbf{59}, 239--244 (1992).
    \href{https://doi.org/10.1007/BF01197321}{DOI: 10.1007/BF01197321}.
\end{thebibliography}
